\documentclass{amsbook}[openany]
\renewcommand\thesection{\thechapter.\arabic{section}}
\usepackage[T1]{fontenc}
\usepackage{amsmath,amssymb}
\usepackage{textcomp}
\usepackage{fancyvrb}
\usepackage{upquote}
\usepackage{hyperref}
\usepackage{booktabs}
\usepackage{longtable}
\usepackage{enumitem}
\newtheorem{thm}{Theorem}[chapter]
\newtheorem{prop}[thm]{Proposition}
\newtheorem{cor}[thm]{Corollary}
\newtheorem{lem}[thm]{Lemma}
\newtheorem{claim}{Claim}

\theoremstyle{definition}
\newtheorem{dfn}{Definition}[chapter]
\newtheorem{example}{Example}[chapter]
\theoremstyle{remark}
\newtheorem{remark}{Remark}[chapter]
\theoremstyle{remark}
\pdfglyphtounicode{visiblespace}{0020}
\pdfgentounicode=1

\DefineVerbatimEnvironment{VNBcode}{Verbatim}{frame=single,framesep=4pt,xleftmargin=8pt,fontsize=\small}
\DefineVerbatimEnvironment{VNBsnip}{Verbatim}{fontsize=\small}

\title{\textbf{VNB Proof Checker}\\[4pt]
       \large User Manual}
\author{F. J. Thayer}
\date{}

\begin{document}
\maketitle
\tableofcontents
\newpage

%% =========================================================
\chapter{Introduction}

The VNB Proof Checker is an interactive proof assistant implemented in
MIT~Scheme.  Its logical foundation is a first-order set theory in the
style of von~Neumann--Bernays (VNB): a single-sorted universe of
\emph{classes}, with \textsc{set} as a distinguished predicate and
membership ($\in$) as the only non-logical primitive.
[[NOTE: References needed. NO HALLUCINATIONS please]]

The proof architecture follows the sequent-calculus approach of the
IMPS~2.0 system~\cite{IMPS}: proofs are \emph{deduction graphs} whose
nodes are sequents and whose edges are inference steps.  Every
inference step is validated by a small, auditable \emph{primitive
inference kernel}; higher-level tactics and a form of conditional
rewrite rules, called macetes in the IMPS system, are built on top of
the inference kernel and cannot introduce unsoundness.
[[Reference to IMPS and an online Portuguese language dictionary
(e.g.\ \url{https://dicionario.acad-ciencias.pt/pesquisa/macete/} has
an entry)]]

This project is, among other things, a proof of concept: a rigorous
mathematical proof system built entirely by AI
[[Note: we can reference documentation for Claude]],
with architectural direction from F.~J.~Thayer.  It is inspired by
ideas from IMPS~\cite{IMPS}.

%% =========================================================
\section{Overview of other proof systems}
\label{sec:other-systems}

\emph{[To be written.  This chapter will give a brief comparative
overview of HOL, Coq, LEAN, and related systems, and will discuss
recent work by Terence Tao on machine-assisted mathematics.  A
literature search is required; no content will be added here without
accurate citations.]}

\section{VNB: primitives and reductions}
\label{sec:enriched}

\emph{[References to VNB in the published literature will be added here
after a literature search.  No content will be added without accurate
citations.]}

Pure VNB set theory is a spare formalism.  The natural numbers, the
reals, functions, and ordinals can all be \emph{constructed} from the
axioms of set theory---but those constructions are long, intricate, and
bear little resemblance to how mathematics is actually written or thought
about.  A proof that $\sin(\pi/2) = 1$ should not need to unfold the
Dedekind-cut definition of the reals.

This system therefore adopts an \emph{enriched} approach.  VNB set
theory supplies the logical foundation, but many mathematical objects
are introduced as \emph{primitive constants} and \emph{primitive
constructors}, governed by axioms rather than by explicit
set-theoretic construction:

\begin{itemize}[label = $\circ$, leftmargin= 0.9cm]
\item \textbf{Number systems} \texttt{NN}, \texttt{ZZ}, \texttt{QQ},
  \texttt{RR}, \texttt{CC} are primitive class constants.  Their
  elements and operations are characterised by ring, field, and order
  axioms, not by Dedekind cuts, Cauchy sequences, or ordinal arithmetic.
\item \textbf{Ordinals} \texttt{ORD} form a primitive proper class with
  its own ordering predicate \texttt{<=\_ORD}, independent of the
  set-membership relation.
\item \textbf{Function spaces} \texttt{fun($A$, $B$)} are primitive
  class constructors.  Functions are not encoded as sets of ordered
  pairs; application \texttt{$f$($x$)} is a primitive operation.
\item \textbf{Mathematical structures} such as groups, vector spaces,
  and Banach spaces are declared with \texttt{def-structure}, which
  installs accessor operations and a typing predicate.
\end{itemize}

In principle every one of these primitives could be reduced to the bare
VNB axioms---ordinals by the von~Neumann construction, reals by Dedekind
cuts, functions by sets of ordered pairs, and so on.  The enriched
approach simply takes these reductions as \emph{existence theorems} that
have been proved (or could be proved) elsewhere, and works with the
objects directly at their natural level of abstraction.  The resulting
proofs are shorter, more readable, and correspond more closely to
informal mathematical practice.

A further departure from standard set theory is the treatment of
\emph{partial application}.  Any syntactic expression \texttt{f(x)}
is a well-formed term, but it may be \emph{undefined} if \texttt{f} is
not a function applicable to \texttt{x}.  The system provides both
strict equality \texttt{a = b} (false if either side is undefined) and
quasi-equality \texttt{a == b} (true when both sides are undefined in
the same way), following the approach of IMPS~\cite{IMPS}.

A consequence of the enriched approach is that membership in a
primitive constant has \emph{indeterminate truth value} in general.
The formula \texttt{x in 7} is syntactically well-formed --- VNB is
single-sorted and \texttt{in} accepts any two class expressions ---
but the system makes no claim about whether it holds. Please note that
this does not mean that the formula is undefined: whether it is true
or not depends on how $7$ is interpreted in VNB. The enriched system
takes $7$ as a primitive number constant and makes no identification
of it with any set-theoretic object.  Membership statements of this
kind are grammatically meaningful but proof-theoretically inert.

%% =========================================================
\section{Architecture}
\label{sec:architecture}

\subsection{Proof architecture}

The proof architecture follows the sequent-calculus approach of the
IMPS~2.0 system~\cite{IMPS}: proofs are \emph{deduction graphs} whose
nodes are sequents and whose edges are inference steps.  Every inference
step is validated by a small, auditable \emph{primitive inference kernel};
higher-level tactics and macetes are built on top of it and cannot
introduce unsoundness.

\subsection{Theory architecture}

In contrast to the \emph{little theories} methodology of
IMPS~\cite{IMPS}---where each mathematical domain inhabits its own
separate mini-theory---VNB works with a \emph{single, growing current
theory}.  Mathematical structures are declared directly in that one theory
using \texttt{def-structure}, and their axioms accumulate there alongside
the base VNB axioms and the primitive number-system axioms.  This reflects
the view of mathematics as one unified body of knowledge rather than a
federation of isolated theories.

The system does permit \emph{minimal} or even \emph{odd-ball theories}
for the cases where a genuinely separate context becomes necessary.
These theories can be created with \texttt{make-theory} and carry
the primitive class constants but \emph{no structural axioms}.  Such
theories are expected to be exceptional; all routine mathematical work
takes place in the single current theory.

\begin{remark}\label{remark:27137-61828-986360}
%
For the potential user interested in \emph{constructivist} mathematics,
this system is not appropriate. These potential users should consider
alternatives such as Coq, which has a very well established library of
mathematics.
%
\end{remark}

\subsection{Design goals}

\begin{itemize}[ label = $\circ$, leftmargin = 1cm]
\item \textbf{Enriched set-theoretic foundation.}  VNB set theory with
  primitive constants for standard mathematical objects; no type
  hierarchy; partial function application with undefinedness.
\item \textbf{Transparent kernel.}  All proof steps reduce to a small
number of
  primitive inferences; the kernel is under 300 lines of Scheme.
\item \textbf{Composable rewriting.}  A macete mechanism (borrowed from
  IMPS) allows theorems to be automatically converted into named rewrite
  strategies that can be sequenced and combined.
\item \textbf{N-ary constructors.}  Ordered tuples \texttt{[$e_1$,
  \ldots, $e_n$]} and Cartesian products \texttt{cartesian($A_1$, \ldots,
  $A_n$)} are first-class, supporting the definition of mathematical
  structures with multiple carriers and operations.
\end{itemize}

\subsection{Why MIT Scheme?}

MIT~Scheme was chosen for several practical reasons: it is small (the
proof kernel is under 300 lines of idiomatic Scheme), easy to learn
and write code, free and readily available on all major platforms, and has
a long track record of correctness and stability.  Since all code is
open source, porting to another language with AI assistance should be
straightforward.

%% =========================================================

%% =========================================================
\section{Quick start}
\label{sec:quickstart-chap}

\subsection{Installation and startup}

\subsubsection{Requirements}

MIT~Scheme~11.2 or later, and for best interactive use one of
GNU/Emacs or XEmacs.  With minor modifications, the native Scheme text
editor Edwin also works.  No other dependencies.

\subsubsection{Directory layout}

\begin{VNBsnip}
~/prover/
  prover                 executable script
  load.scm               loader (sources all files in order)
  expressions.scm        expression language
  wff.scm                <wff> record type, wff-child, wff-in-theory, wff-equiv?
  sequents.scm           sequents over <wff>; context operations
  deduction-graphs.scm   proof graph data structures
  primitive-inferences.scm  inference kernel
  macetes.scm            macete mechanism
  theory.scm             theory management + VNB base axioms
  structures.scm         mathematical structure definitions
  number-systems.scm     NN, ZZ, QQ, RR, CC constants and axioms
  contexts.scm           local context management, make-wff, unravel
  arith-eval.scm         ground arithmetic decision procedure
  proof-commands.scm     high-level proof commands
  interactive.scm        interactive short-form commands (sp, di, ai, ...)
  vnb.el                 GNU Emacs / XEmacs interface
\end{VNBsnip}

\subsection{Starting the checker}

\begin{VNBsnip}
~/prover/prover              # interactive REPL
~/prover/prover proof.scm    # run a proof script and exit
~/prover/prover -i proof.scm # run a script, then stay interactive
\end{VNBsnip}

The interactive session is a standard MIT~Scheme REPL with the proof
checker pre-loaded.
\subsection{The Emacs interface}
\label{sec:emacs}

The file \texttt{vnb.el} provides a GNU~Emacs and XEmacs interface for
the proof checker.  It requires no external packages beyond
\texttt{comint}, which ships with every Emacs distribution.

\subsubsection{Installation}

Add the following to your Emacs init file (\texttt{\textasciitilde/.emacs}
or \texttt{\textasciitilde/.xemacs/init.el}):

\begin{verbatim}
(add-to-list 'load-path "~/prover")
(require 'vnb)
\end{verbatim}

\noindent
\textbf{Do not byte-compile \texttt{vnb.el}.}  GNU~Emacs and XEmacs use
incompatible byte-code formats; loading the source file interpreted is
the correct approach for both.

\subsubsection{Starting the interface}

\begin{VNBsnip}
M-x vnb
\end{VNBsnip}

This starts the prover subprocess (if not already running), and opens a
two-buffer layout:
\begin{itemize}[label=\textendash, leftmargin=1cm]
\item \textbf{Left window --- \texttt{*VNB*} REPL.}  A \texttt{comint}
  buffer running MIT~Scheme with the proof checker pre-loaded.  Type any
  Scheme expression and press \texttt{RET} to evaluate it.
\item \textbf{Right window --- \texttt{*VNB State*}.}  A read-only
  display buffer that is updated automatically after every command that
  calls \texttt{(show)}.  It always shows the current proof state: open
  goals, assumptions, and focus.
\end{itemize}

The two buffers are kept synchronised by a comint output filter that
watches for the sentinel lines \texttt{;;VNB-STATE-BEGIN} /
\texttt{;;VNB-STATE-END} emitted by \texttt{(show)} and routes the
enclosed text to the state buffer.

\subsubsection{Key bindings in the \texttt{*VNB*} REPL buffer}

\begin{description}
\item[\texttt{RET}] Submit expression to the prover (standard comint).
\item[\texttt{C-c C-s}] Send \texttt{(show)} to refresh the state display.
\item[\texttt{C-c C-p}] Switch to the \texttt{*VNB State*} buffer.
\end{description}

\subsubsection{Key bindings in the \texttt{*VNB State*} buffer}

\begin{description}
\item[\texttt{q}] Return to the \texttt{*VNB*} REPL buffer.
\item[\texttt{C-c C-s}] Send \texttt{(show)} to refresh the state display.
\end{description}

\subsubsection{The \texttt{*VNB Commands*} scratch buffer}

\begin{VNBsnip}
M-x vnb-command-buffer
\end{VNBsnip}

Opens a \texttt{*VNB Commands*} buffer in \texttt{vnb-command-mode}.
This buffer works exactly like the Emacs \texttt{*scratch*} buffer (which
uses \texttt{lisp-interaction-mode}), except that \texttt{C-j} sends
expressions to the \emph{VNB prover} instead of evaluating them as Emacs
Lisp.

Write proof commands as S-expressions in the buffer.  Position point
immediately after any complete expression and press \texttt{C-j}:
the expression is sent to the running VNB process, and the return value
is inserted on a new line below point.  The buffer accumulates a
readable session transcript.

\begin{VNBcode}
(sp (make-wff-from-string "n in NN and m in NN implies m in NN and n in NN"))
;;VNB-STATE-BEGIN
1 open goal(s).
Focus:
  [0] ()  =>  (((n in nn) and (m in nn)) implies ((m in nn) and (n in nn)))
;;VNB-STATE-END

(di)       C-j
;;VNB-STATE-BEGIN
2 open goal(s).
Focus:
  [1] ((n in nn) and (m in nn))  =>  ((m in nn) and (n in nn))
Other open goals:
  [0] ()  =>  (((n in nn) and (m in nn)) implies ((m in nn) and (n in nn)))
;;VNB-STATE-END
\end{VNBcode}

\paragraph{Key bindings in \texttt{vnb-command-mode}.}

\begin{description}
\item[\texttt{C-j}] Send S-expression before point to VNB; insert result below.
\item[\texttt{C-c C-c}] Send entire buffer to VNB; append result.
\item[\texttt{TAB}] Complete VNB command name at point.
\item[\texttt{C-c C-d}] Describe VNB command: show signature and documentation.
\item[\texttt{C-c C-s}] Refresh \texttt{*VNB State*} proof-state buffer.
\end{description}

\subsubsection{Command reference and completion}

\texttt{vnb.el} maintains the list \texttt{vnb-commands} of all proof
commands, each with a usage signature and a one-line description.  The
derived alist \texttt{vnb-commands-alist} maps short names to signatures
and is used by \texttt{TAB} completion in \texttt{vnb-command-mode}.

\begin{VNBsnip}
M-x vnb-describe-command   RET   di   RET
; Echo area: (di)
;   Direct inference: decompose current goal by its top connective. ...
\end{VNBsnip}

\subsubsection{Programmatic evaluation: \texttt{vnb-eval-string}}

For scripting or integration with other Emacs packages:

\begin{VNBsnip}
; str is a Scheme expression as a string
(vnb-eval-string str)           ; => Scheme return value as string
(vnb-eval-string str timeout)   ; => same, with explicit timeout (seconds)
\end{VNBsnip}

\texttt{(vnb-eval-string STR \&optional TIMEOUT)} sends the string
\texttt{STR} to the running prover, waits up to \texttt{TIMEOUT} seconds
(default 30) for the next REPL prompt, and returns the Scheme return
value as an Emacs string.  It signals an error if the prover is not
running.

%% =========================================================
\subsection{A complete proof}
\label{sec:quickstart}

After installing the Emacs interface (Section~\ref{sec:emacs}), launch
the prover with \texttt{M-x vnb}.  Two buffers appear: \texttt{*VNB*}
(the REPL) and \texttt{*VNB State*} (the live proof state).  Type the
following five commands at the \texttt{*VNB*} prompt, pressing
\texttt{C-j} (or RET) after each:

\begin{VNBsnip}
(sp (make-wff-from-string "n in NN and m in NN implies m in NN and n in NN"))
(di)                                   ; assume both, prove (m in nn) and (n in nn)
(ai "n in NN and m in NN")            ; split into (n in nn) and (m in nn)
(di)                                   ; split goal: prove (m in nn), then (n in nn)
(ass) (ass)                            ; both are in the assumption list
\end{VNBsnip}

The \texttt{*VNB State*} buffer updates after each command.  The final
state shows \texttt{Proof complete.}

\paragraph{What just happened.}
\texttt{(make-wff-from-string \ldots)} parses a string in the VNB string
syntax and wraps the result as a typed \texttt{wff} in the current
theory.  \texttt{(sp \ldots)} starts a new proof.  \texttt{(di)} is
direct-inference: it decomposes the goal by its outermost connective.
\texttt{(ai \ldots)} is antecedent-inference: it decomposes a named
assumption; the argument is a formula string (or a raw S-expression for
low-level use) that identifies the assumption by alpha-equivalence.
\texttt{(ass)} discharges a goal that already appears in the assumption
list.

The same proof at the Scheme level (no short forms) reads:
\begin{VNBsnip}
(let* ((wic (make-wff-from-string
              "n in NN and m in NN implies m in NN and n in NN"))
       (ps  (start-proof wic))
       (ps  (cmd-direct-inference ps))
       (ps  (cmd-antecedent-inference ps (parse-string "n in NN and m in NN")))
       (ps  (cmd-direct-inference ps))
       (ps  (cmd-assumption ps))
       (ps  (cmd-assumption ps)))
  (print-proof-state ps))
\end{VNBsnip}
The interactive short forms in \texttt{interactive.scm} are simple
wrappers: each mutates a global \texttt{*ps*} and calls \texttt{(show)}
to reprint the state.  All examples in this manual use the short forms;
the long form is available when scripting non-interactively.

\paragraph{What \texttt{make-wff} checks.}
\texttt{make-wff} (and \texttt{make-wff-from-string}, which calls it)
performs a position-typed walk of its argument and signals an error
rather than wrap a malformed expression.  The checks are:
\begin{itemize}[label=$\ast$, leftmargin=1cm]
\item \textbf{Numbers in wff positions} are rejected:
  \texttt{(make-wff-from-string "1 and 35")} $\Rightarrow$ ``number in wff
  position 1''.
\item \textbf{Term-forming heads in wff positions} are rejected:
  \texttt{(make-wff-from-string "choice(A)")} as a top-level wff fails because
  \texttt{choice} produces a term, not a wff.
\item \textbf{Wff-forming heads in term positions} are rejected:
  \texttt{(make-wff-from-string "(A and B) in C")} fails because \texttt{A and B}
  is a wff and the first argument of \texttt{in} must be a term.
\item \textbf{Bare symbols in wff positions} are rejected.  Only
  \texttt{TRUTH} and \texttt{FALSEHOOD} are valid bare-symbol wffs;
  anything else (e.g.\ \texttt{A implies B}) is an error.
\item \textbf{Free-symbol role conflicts}: a symbol that appears free
  in both wff and term positions in the same formula is rejected.
\item \textbf{Wff-position binders}: a binder whose body uses the
  bound variable in wff position is rejected within its scope.  This
  catches \texttt{forall([p], p implies p)}: \texttt{p} is bound as a
  class variable but the body uses it as a proposition.
\item \textbf{Bound vs.\ free same-name use} is \emph{accepted with a
  warning}.  In \texttt{forall([a], a in nn) implies a in zz} the
  inner \texttt{a} is bound and the outer \texttt{a} is free.  The
  validator emits a \texttt{;; warning} but does not reject.
\item \textbf{Arities and bound-variable shapes} are checked: every
  connective and quantifier gets the right number of arguments;
  \texttt{FORALL}, \texttt{FORSOME}, \texttt{IOTA},
  \texttt{VNB-LAMBDA}, and \texttt{SEP} bound variables must be
  symbols.
\end{itemize}
%
\begin{remark}\label{remark:27135-34761-138906}
%
The set-theoretic \texttt{power(A)} (power set) and arithmetic
\texttt{x \^{} n} (exponentiation) share the same underlying symbol because
MIT~Scheme reads identifiers case-insensitively; \texttt{make-wff}
distinguishes them by arity (1 vs.\ 2 arguments).  All downstream
operations (\texttt{free-vars}, \texttt{subst-free},
\texttt{alpha-equiv?}, macete pattern matching) honour the same
arity-based dispatch.

%
\end{remark}
%% =========================================================

\chapter{Formal structure}
\label{sec:language}

The language has two syntactic layers.

The \emph{internal layer} represents every expression as a plain
Lisp S-expression.  This representation is time-tested and deliberately
simple: it requires no parser, supports straightforward recursive
processing, and is directly inspectable and constructible at the Scheme
REPL.  All prover machinery---inference rules, substitution, pattern
matching, macetes, arithmetic evaluation---operates on this layer.
Users are not required to engage with it directly; however, in the
tradition of early Lisp, they are free to do so.  Constructing or
manipulating S-expressions by hand, loading axioms as quoted lists,
or writing Scheme code that builds formulas programmatically are all
legitimate and supported uses.

The \emph{string syntax} is the surface notation accepted by
\texttt{make-wff-from-string} and emitted by \texttt{wff->string}
/ \texttt{pp}: binary operators are written infix, multi-argument
terms use functional notation, quantifiers use the binding-list form,
and set expressions use brace notation.
Section~\ref{sec:parser-printer} describes both layers in full.
The examples in this manual use string syntax for user-facing
formulas; internal S-expression forms appear only in theory and
structure definitions.

MIT~Scheme reads symbols case-insensitively (folding to lower case),
so \texttt{FORALL}, \texttt{Forall}, and \texttt{forall} are the same
symbol at the Scheme level.

As should be the case in any formal system, we distingish between
syntax and semantics. In what follows, we distinguish two distinct
types of syntactic categories, each with its own semantics: formulas
and terms. These syntactic categories can be viewed in either the
internal syntax or the string syntax. In the presentation below, we
adhere to the string syntax.


\section{Formulas}

Semantically, a formula is a syntactic expression that can have only
two values: TRUTH or FALSEHOOD. 
%
\begin{remark}\label{remark:27137-52766-301830}
%
Note that ``UNDEFINED'' is not a possible value for a
formula. However, the value of any particular formula may not be
known. In fact, since we are silent about how elements in our system
may be interpreted in a pure von Neumann Bernays theorem, different
interpretations may yield different values.
%
\end{remark}

Formulas can be produced from other ones by basic logical operators,
including quantification over variables that range over objects.
%
In the following table, the symbols $p$, $q$ are not variables that
range over formulas, but rather placeholders that can instantiated
with well-formed formulas.

\begin{itemize}[label=\textendash, leftmargin=1cm]
\item \texttt{TRUTH} propositional truth $\top$
\item \texttt{FALSEHOOD} propositional falsity $\bot$
\item \texttt{not($p$)} negation $\lnot p$
\item \texttt{$p$ and $q$} conjunction $p \land q$
\item \texttt{$p$ or $q$} disjunction $p \lor q$
\item \texttt{$p$ implies $q$} implication $p \Rightarrow q$
\item \texttt{$p$ iff $q$} biconditional $p \Leftrightarrow q$
\item \texttt{forall([$b_1$, \ldots, $b_n$], $p$)} universal quantification; see Section~\ref{sec:multiquant}
\item \texttt{forsome([$b_1$, \ldots, $b_n$], $p$)} existential quantification; see Section~\ref{sec:multiquant}
\item \texttt{$a$ = $b$} strict equality: both defined and equal
\item \texttt{$a$ == $b$} quasi-equality: both defined and equal, or both undefined
\item \texttt{$a$ in $b$} membership $a \in b$
\item \texttt{$A$ subset $B$} subclass $A \subseteq B$; see Section~\ref{sec:class-set-discipline}
\end{itemize}

\noindent
The quantifiers \texttt{forall} and \texttt{forsome} range over \emph{all
classes} (the system is single-sorted).  There is no restriction to sets.

\subsection{Quantification}
\label{sec:multiquant}

Both \texttt{forall} and \texttt{forsome} require a bracket-enclosed
\emph{binding list}, even for a single variable:

\begin{itemize}[label = $\circ$, leftmargin=1cm]
\item \texttt{forall([binding, ...], body)}

\item \texttt{forsome([binding, ...], body)}
\end{itemize}

\noindent
\texttt{forall(x, x = x)} is a syntax error; the correct form is
\texttt{forall([x], x = x)}.

\noindent
Each \emph{binding} is one of:

\begin{description}
\item[\texttt{$x$}] unrestricted variable ranging over all classes
\item[\texttt{$x$ in $A$}] variable $x$ restricted to class $A$
\item[\texttt{[$n_1$, \ldots, $n_k$] in $A$}] tuple destructuring: fresh $t \in A$, components bound to \texttt{nth(1,~$t$)}, \ldots, \texttt{nth($k$,~$t$)}
\end{description}

\noindent
Bindings are processed left to right, outermost first.  Thus in the
\emph{internal syntax} they expand to fully nested single-variable form
as follows:

\begin{VNBsnip}
forall([x, y in NN, z], y >= 0)
; => (FORALL x (FORALL y (IMPLIES (IN y NN) (FORALL z (>= y 0)))))

forall([[E, +, *] in METRIC], IS-METRIC(E, +, *))
; => (FORALL t_0 (IMPLIES (IN t_0 METRIC)
;                         (IS-METRIC (NTH 1 t_0) (NTH 2 t_0) (NTH 3 t_0))))
\end{VNBsnip}

\paragraph{Display.}
The printer reconstructs binding-list form from the internal nested
representation.  Consecutive same-type quantifiers are collapsed into
a single binding list; a change of quantifier type stops the collapsing.
For example, \texttt{forall([x in NN], forall([y], p))} prints as
\texttt{forall([x in NN, y], p)}.

\paragraph{Variable shadowing.}
Bound variable names may be reused within nested quantifiers.
The inner binding shadows the outer; \texttt{fresh-var} in
\texttt{subst-free} ensures substitution never captures a free variable.
Such formulas are valid but confusing; users should avoid them.

\subsection{Equality and quasi-equality}

Terms may be \emph{undefined} (e.g.\ an application \texttt{(f x)} when
\texttt{f} is not a function applicable to \texttt{x}).  The system
provides two equality predicates:

\paragraph{Strict equality \texttt{$a$ = $b$}.}
True iff both \texttt{a} and \texttt{b} are defined and equal.
False if either is undefined.  In particular:
\begin{itemize}[label=$-$, leftmargin=1cm]
\item \texttt{george(s) = 75} is \emph{false} when \texttt{george}
  is not a function.
\item \texttt{george(s) = george(s)} is \emph{false} for the same
  reason.
\item A term $t$ is defined iff $t = t$:
  $\mathit{DEF}(t) \;\equiv\; (t = t)$.
\end{itemize}

\paragraph{Quasi-equality \texttt{$a$ == $b$}.}
True iff both \texttt{a} and \texttt{b} are defined and equal, \emph{or}
both are undefined.  Key properties:
\begin{itemize}[label=$-$, leftmargin=1cm]
\item \texttt{t == t} is \emph{always true}, for any term $t$,
  defined or not.
\item \texttt{george(s) == george(s)} is \emph{true}.
\item \texttt{$a$ = $b$} implies \texttt{$a$ == $b$}, but not vice versa.
\item \texttt{$a$ == $b$} and $\mathit{DEF}(a)$ together imply
  \texttt{$a$ = $b$}.
\end{itemize}

\noindent
Undefinedness propagates through all operations: if \texttt{george(s)}
is undefined then \texttt{george(s) + 1} is also undefined,
and \texttt{george(s) + 1 = 2} is false.

Variables and class constants (\texttt{ORD}, \texttt{RR},
\texttt{EMPTY-SET}, \ldots) are always defined.  The application
\texttt{f(x)} is defined when \texttt{f in fun(A, B)} and
\texttt{x in A} are provable for some $A$, $B$.

\paragraph{Reflexivity in proofs.}
\texttt{cmd-quasi-reflexivity} closes any goal of the form
\texttt{a == a} unconditionally.
\texttt{cmd-reflexivity} closes \texttt{a = a} only when \texttt{a} is
demonstrably defined.

Our rendition of von Neumann/Bernays, which we label VNB, recognizes
three kinds of object.  A \emph{set} is an element of the class
$\mathit{SET}$.  A \emph{proper class} is another kind of collection
that can have elements, but they themselves are \emph{not} elements of
anything.  A \emph{Functoid} is an operation that may be defined on
all elements of a class.  All the constructors we list below are
functoids.

\texttt{SET} is a primitive class constant.  The formula
\texttt{$a$ in SET} expresses ``$a$ is a set''.
There is no separate predicate form \texttt{SET($a$)}.
%
The key axiom distinguishing VNB from typed systems:
\[
  (a \in b) \;\Longrightarrow\; a \in \mathit{SET}
\]
Only sets can be members of anything: membership forces the \emph{left}
argument to be a set.  The right argument---the containing class---may be
a proper class.  (For example, $3 \in \mathit{ORD}$ is well-formed even
though the class of all ordinals is a proper class.)

\label{sec:class-set-discipline}


\begin{remark}\label{remark:27137-54738-896234}
%
The introduction of functoids is a convenience, but not an essential
departure from the von Neuman/Bernays framework. In the classical
set-theoretic reduction of functions as ordered pairs, we can view
functoids as \emph{classes} of ordered pairs with the function
property.  The formal language quantifies over sets and classes.
However, we don't specify where ``functoids'' live in the VNB
universe, so statements of the form ``all functoids $F$
satisfy\ldots'' is metalanguage.  Functoids are introduced either from
the predefined list or via \texttt{def-functoid}
(Section~\ref{sec:def-functoid}).

%
\end{remark}

\texttt{apply-function} and \texttt{apply-functoid} are the basic term
constructors and we use the standard abbreviated form for function
application as follows:
%
\begin{itemize}[label=$\cdot$, leftmargin=1cm]
%
\item If $f$ is a function, $f(x_1, \ldots, x_n)$ means $\texttt{apply-function}(f,x_1, \ldots, x_n)$

%

\item If $f$ is a functoid, $f(x_1, \ldots, x_n)$ means
$\texttt{apply-functoid}(f,x_1, \ldots, x_n)$
%
\end{itemize}



%

%

%
The subset relation
\[
  A \subseteq B \;\equiv\; \forall x.\; x \in A \Rightarrow x \in B
\]
is a primitive binary predicate written \texttt{A subset B}
(\texttt{subset-def}).  An immediate corollary of
\texttt{membership-implies-sethood} is that \emph{every class is a
subclass of $\mathit{SET}$} (axiom \texttt{subset-set}): only sets can
appear on the left of~$\in$.

The constructors below are functoids whose stated domains are elements
of $\mathit{SET}$.  Class comprehension $\{x \mid p\}$ is the sole
exception: It is an \emph{schema}, which yields an infinite number of
instances, one for each syntactically correct formula $p$, and each
one with $0$ arguments.

\begin{description}

\item[\texttt{union($A_1$,\ldots,$A_n$)}, $A_i \in \mathit{SET}$]
  $x \in C$ iff $x \in A_i$ for some $i$.

\item[\texttt{intersection($A_1$,\ldots,$A_n$)}, $A_i \in \mathit{SET}$]
  $x \in C$ iff $x \in A_i$ for all $i$
  (requires separation, pending; Section~\ref{sec:axioms}).

\item[\texttt{complement-in($A$, $B$)}, $A \in \mathit{SET}$]
  Relative complement $A \setminus B$: $x \in C$ iff $x \in A$ and
  $x \notin B$ (separation axiom, Section~\ref{sec:axioms}).
  $B$ need not be a set.

\item[\texttt{cartesian($A_1$,\ldots,$A_n$)}, $A_i \in \mathit{SET}$]
  $x \in C$ iff $x = [a_1,\ldots,a_n]$ with $a_i \in A_i$.

\item[\texttt{tuples($A$)}, $A \in \mathit{SET}$]
  $x \in C$ iff $x = [a_1,\ldots,a_n]$ with all $a_i \in A$,
  for some $n \ge 0$.

\item[\texttt{fun($A$, $B$)}, $A, B \in \mathit{SET}$]
  $x \in C$ iff $x$ is a total function with domain $A$ and values in $B$
  (see Section~\ref{sec:funapp}).
  $\mathit{FUN}(A,B) \subseteq \mathit{FUN}(A,C)$ whenever $B \subseteq C$.

\item[\texttt{\{$x$ in $A$: $p$\}} / \texttt{\{$x$ in $A$ | $p$\}}, $A \in \mathit{SET}$]
  Separation: $x \in C$ iff $x \in A$ and $p[x]$ holds
  (axiom pending; Section~\ref{sec:axioms}).

\item[\texttt{\{$x$ | $p$\}}]
  Class comprehension: $x \in C$ iff $p[x]$ holds.  Proper class in
  general unless provably bounded.

\item[\texttt{power($A$)}, $A \in \mathit{SET}$]
  $x \in C$ iff $x \in \mathit{SET}$ and $x \subseteq A$.

\end{description}

\begin{remark}
\texttt{union($A_1$, \ldots, $A_n$)} takes a finite listed union.
The \emph{big-union} $\bigcup A = \{x \mid \exists a \in A.\; x \in a\}$
of an arbitrary class $A$ of sets is a separate constructor
\texttt{UNION-SET}, currently pending (Section~\ref{sec:axioms}).
\end{remark}

\section{Term constructors}
\label{sec:choice-iota-lambda}

\emph{Term constructors} build syntactic objects whose intended
reprsentation are objects in the VNB universe.

\begin{description}
\item[\texttt{[$e_1$, \ldots, $e_n$]}] ordered $n$-tuple $(e_1, \ldots, e_n)$; always a set
\item[\texttt{nth($k$, $e$)}] $k$-th component of tuple $e$ (1-indexed); $k$ may be any term;
  \texttt{nth($k$, [$e_1$,\ldots,$e_n$])} reduces to $e_k$ via
  \texttt{cmd-nth-reduce} (Section~\ref{sec:commands})
\item[\texttt{length($L$)}] number of elements in tuple $L$; a natural number:
  \begin{itemize}[label=$-$, leftmargin=1cm]
  \item \texttt{length-of-empty}: $\texttt{length}([\,]) = 0$
  \item \texttt{length-in-nn}: $L \in \mathit{TUPLES}(\mathit{SET}) \Rightarrow
    \texttt{length}(L) \in \mathit{NN}$
  \item \texttt{nth-in-range}: $L \in \mathit{TUPLES}(A),\; i \in \mathit{NN},\;
    1 \le i \le \texttt{length}(L) \Rightarrow \texttt{nth}(i,L) \in A$
  \end{itemize}
  A full recursive characterisation requires a \texttt{prepend} constructor
  and the \texttt{tuples-induction} schema, both pending
  (Section~\ref{sec:pending-axioms}).
\item[\texttt{EMPTY-SET}] the empty set $\emptyset$; always a set; $= \texttt{make-set}([\,])$
\item[\texttt{\{$a$\}}] singleton $\{a\}$; always a set; $= \texttt{make-set}([$a$])$
\item[\texttt{make-set($L$)}] unordered finite set from tuple $L \in \mathit{TUPLES}(A)$,
  $A \in \mathit{SET}$; always a set:
  \begin{itemize}[label=$-$, leftmargin=1cm]
  \item \texttt{make-set-membership}: $x \in \texttt{make-set}(L)
    \Longleftrightarrow \exists\,i \in \mathit{NN}.\; \texttt{nth}(i,L) = x$
  \item \texttt{make-set-sethood}: $L \in \mathit{TUPLES}(A),\; A \in \mathit{SET}
    \Rightarrow \texttt{make-set}(L) \in \mathit{SET}$
  \item \texttt{make-set-empty}: $\texttt{make-set}([\,]) = \emptyset$
  \end{itemize}
  The notation $\{a,b,c,\ldots\}$ is sugar for $\texttt{make-set}([a,b,c,\ldots])$;
  repetitions are harmless.
\item[\texttt{choice($A$)}] primitive choice operator: an element of nonempty class~$A$
  (undefined when $A$ is empty).  Axiom \texttt{choice-axiom}:
  \[
    (\exists x.\; x \in A) \;\Longrightarrow\; \texttt{choice}(A) \in A
  \]
  This is the Bourbaki-style global choice: a single canonical operator
  defined uniformly on all nonempty classes, not just sets.
\item[\texttt{iota($x$, $p$)}] definite description: the unique object satisfying~$p$;
  a binder.  Defined only when $\exists! x.\; p(x)$ is provable:
  \[
    \texttt{iota}(x, p) \in A \quad\text{requires a proof of}\quad \exists! x.\; p(x)
  \]
  A planned \texttt{define-object} form will wrap this pattern:
\begin{VNBsnip}
(define-object sqtwo (IOTA x (= (* x x) 2)))
; requires a proved theorem: (FORSOME-UNIQUE x (= (* x x) 2))
; installs sqtwo as a constant with (= (* sqtwo sqtwo) 2)
\end{VNBsnip}
\item[\texttt{vnb-lambda($x$, $\mathit{body}$)}] anonymous function $x \mapsto \mathit{body}$;
  a binder.  The name avoids clashing with Scheme's \texttt{lambda}.
  Typing axiom (pending):
  \[
    \bigl(\forall x \in A.\; \mathit{body}(x) \in B\bigr)
    \;\Longrightarrow\;
    \texttt{vnb-lambda}(x, \mathit{body}) \in \mathit{FUN}(A,B)
  \]
  $\beta$-reduction (\texttt{cmd-vnb-lambda-beta}, short form \texttt{(beta)}):
  \[
    \texttt{vnb-lambda}(x, \mathit{body})(v) \;=\; \mathit{body}[x \mapsto v]
  \]
\end{description}

\begin{remark}
Tuples $[e_1,\ldots,e_n]$ are primitive here.  One could in principle reduce
them to iterated pairs (Kuratowski style)[[REFERENCE NEEDED]], but we have no interest in doing so:
$[\ ]$ is perfectly well-defined as a term constructor of arity zero, and that
is all we need.
\end{remark}

\begin{remark}
The index $k$ in $\texttt{nth}(k, e)$ may be any term (not only a literal
integer).  The primitive inference \texttt{cmd-nth-reduce} does require a
literal integer to reduce $\texttt{nth}(k, [e_1,\ldots,e_n])$ to $e_k$, but
$\texttt{nth}(i, L)$ with a variable index $i$ is a perfectly valid term and
appears in \texttt{make-set-membership} and \texttt{nth-in-range}.
\end{remark}

\begin{remark}
\texttt{choice}, \texttt{iota}, and \texttt{vnb-lambda} all participate
correctly in \texttt{free-vars}, \texttt{subst-free}, \texttt{alpha-equiv?},
and macete pattern-matching.
\end{remark}

\section{Function application}
\label{sec:funapp}

Function application is written as a compound term: \texttt{$f$($x$)} for
any object $f$ and any argument $x$.  Examples:
\begin{VNBsnip}
a + b               ; addition applied to a, b
NORM(E, v)          ; NORM of Banach space E applied to vector v
GEORGE(s)           ; syntactically valid, but may be undefined
\end{VNBsnip}

\textbf{FUN($A$, $B$) with $A, B \in \mathit{SET}$.}
$\mathit{FUN}(A, B)$ is a primitive class constructor: its members are
the total functions from $A$ to $B$.  Since $A, B \in \mathit{SET}$,
$\mathit{FUN}(A,B)$ is itself a set (\texttt{fun-set-iff}), and each
$f \in \mathit{FUN}(A,B)$ is a set (\texttt{membership-implies-sethood}).

\begin{remark}
The soundness of $\mathit{FUN}(A,B)$ is confirmed by the standard
set-theoretic construction: a function $f : A \to B$ is a set of pairs
$(a, b) \in A \times B$ that is functional (each $a$ has at most one $b$)
and total (every $a \in A$ appears).  Since $A \times B$ is a set, so
is any subset of it, making $f$ a set.
\end{remark}


\begin{remark}
Category theory motivates a natural extension of $\mathit{FUN}$.  A
functor between large categories---such as the category of all groups or
all topological spaces---maps a proper class of objects to a proper class
of objects.  Since $\mathit{FUN}$ requires a set-valued domain, functors
fall outside $\mathit{FUN}$ entirely.

The planned treatment introduces \emph{operators} as a third kind of
entity alongside sets and classes.  There is no class $\mathit{OPER}(A,B)$
(no such collection can exist without violating the set/class dichotomy),
but there will be a predicate
\[
  \mathit{ISOPER}(A,\, B,\, f)
\]
asserting that $f$ behaves as a total function from $A$ to $B$ even when
$A$ or $B$ is a proper class.  An accompanying term constructor
$\texttt{apply-operator}(f, x)$ then gives the value of $f$ at $x$, for
any $x \in A$, under the hypothesis $\mathit{ISOPER}(A,B,f)$.

Operators are not members of anything---they do not appear on the left of
$\in$---so they are consistent with
\texttt{membership-implies-sethood}.  The three-level ontology (sets,
classes, operators) keeps the VNB foundation intact while accommodating
the full generality of categorical constructions.
\end{remark}

%
The base theory installs the following axioms:

\begin{description}
\item[\texttt{fun-apply-def}]
  $f \in \mathit{FUN}(A,B) \land x \in A \;\Rightarrow\; f(x) = f(x)$
\item[\texttt{fun-apply-in}]
  $f \in \mathit{FUN}(A,B) \land x \in A \;\Rightarrow\; f(x) \in B$
\item[\texttt{fun-extensionality}]
  $f,g \in \mathit{FUN}(A,B) \land (\forall x \in A.\;f(x) = g(x)) \;\Rightarrow\; f = g$
\item[\texttt{fun-set-iff}]
  $\mathit{FUN}(A,B) \in \mathit{SET} \;\Leftrightarrow\; (A \in \mathit{SET}) \land (B \in \mathit{SET})$
\item[\texttt{fun-monotone}]
  $B \subseteq C \;\Rightarrow\; \mathit{FUN}(A,B) \subseteq \mathit{FUN}(A,C)$
\item[\texttt{cartesian-set-iff}]
  $A \times B \in \mathit{SET} \;\Leftrightarrow\; (A \in \mathit{SET}) \land (B \in \mathit{SET})$
\item[\texttt{tuples-sethood}]
  $A \in \mathit{SET} \;\Rightarrow\; \mathit{TUPLES}(A) \in \mathit{SET}$
\end{description}

\begin{remark}
The constructor $\mathit{TUPLES}(A)$ is the class of all finite sequences
$[a_1, \ldots, a_n]$ with each $a_i \in A$, for any $n \ge 0$ (including
the empty sequence $[\,]$).  It fills a gap that $\mathit{CARTESIAN}$ does
not cover: $\mathit{CARTESIAN}(A_1, \ldots, A_n)$ fixes the arity $n$ at
the time the expression is written, whereas $\mathit{TUPLES}(A)$ lets $n$
range over all natural numbers.

The canonical use is \emph{n-ary functions on $A$}: an element of
$\mathit{FUN}(\mathit{TUPLES}(A), B)$ accepts a finite sequence of elements
of $A$ of any length and returns a value in $B$.  For example, addition on
$\mathbb{C}$ (as a variadic operation) would be an element of
$\mathit{FUN}(\mathit{TUPLES}(\mathit{CC}), \mathit{CC})$.
Since $\mathit{TUPLES}(A) \in \mathit{SET}$ when $A \in \mathit{SET}$,
this is a perfectly ordinary instance of $\mathit{FUN}(A', B)$ with
$A' = \mathit{TUPLES}(A)$.
\end{remark}

\section{Constructor examples}
\label{sec:constructor-examples}

A quick reference showing every functoid and term constructor in concrete
VNB notation.  Each example is a valid expression; the comment gives the
key membership or typing fact it participates in.

\paragraph{Functoids.}
\begin{VNBsnip}
union(NN, ZZ, QQ)
; x in union(NN,ZZ,QQ)  iff  x in NN or x in ZZ or x in QQ

intersection(QQ, ORD)
; x in intersection(QQ,ORD)  iff  x in QQ and x in ORD

complement-in(ZZ, NN)
; x in complement-in(ZZ,NN)  iff  x in ZZ and not (x in NN)

cartesian(NN, ZZ)
; [3, -1] in cartesian(NN,ZZ)   requires  3 in NN  and  -1 in ZZ

tuples(NN)
; [3, 7, 2] in tuples(NN)   requires  3,7,2 in NN   (any length n >= 0)
; []        in tuples(NN)   (empty sequence; always a member)

fun(NN, ZZ)
; f in fun(NN,ZZ)  means f is total on NN with values in ZZ

{n in NN: n <= 100}
; separation: x in {n in NN: n<=100}  iff  x in NN and x <= 100

{x | x in NN and x in ORD}
; class comprehension; proper class in general unless provably bounded

power({1, 2})
; x in power({1,2})  iff  x in SET  and  forall z in x. z in {1,2}

\end{VNBsnip}

\paragraph{Term constructors.}
\begin{VNBsnip}
EMPTY-SET
; always a set; no members

{a}
; singleton; a in {a}

{1, 2}        ; = make-set([1, 2]); always a set
{1, 2, 3, 4}  ; = make-set([1, 2, 3, 4]); always a set; repetitions ignored

[a, b]        ; ordered pair; [a,b] in cartesian(A,B) when a in A, b in B
[a, b, c]     ; ordered triple; similarly for longer tuples

nth(1, [a, b, c])   ; = a    (1-indexed)
nth(2, [a, b, c])   ; = b

choice(NN)
; (exists x. x in NN) implies choice(NN) in NN

iota(x, x * x = 4 and x in NN)
; = 2, provided (forsome-unique x. x*x=4 and x in NN) has been proved

vnb-lambda(n, n + 1)
; anonymous function  n |-> n+1
; vnb-lambda(n,n+1) in fun(NN,NN)  when  forall n in NN. n+1 in NN  is proved
; (vnb-lambda(n, n+1))(3)  beta-reduces to  4
\end{VNBsnip}

%% =========================================================
\section{Basic arithmetic}
\label{sec:basic-arith}

\subsection{Complex numbers} \label{sec:nums}

The file \texttt{number-systems.scm} declares five number systems as
class constants and installs their arithmetic axioms in the current
theory.  All declarations take effect at load time.

\subsubsection{The five constants}

\begin{description}
\item[\texttt{NN}] ($\mathbb{N}^+$) Positive integers $\{1, 2, 3, \ldots\}$ --- no zero
\item[\texttt{ZZplus}] ($\mathbb{Z}_{\geq 0}$) Non-negative integers $\{0, 1, 2, \ldots\}$
\item[\texttt{ZZ}] ($\mathbb{Z}$) Integers
\item[\texttt{QQ}] ($\mathbb{Q}$) Rationals
\item[\texttt{RR}] ($\mathbb{R}$) Reals
\item[\texttt{CC}] ($\mathbb{C}$) Complex numbers
\end{description}

Sethood axioms (\texttt{nn-is-set} through \texttt{cc-is-set},
\texttt{zzplus-is-set}) and the inclusion chain
\[
  \mathbb{N}^+ \subseteq \mathbb{Z}_{\geq 0} \subseteq \mathbb{Z}
                \subseteq \mathbb{Q} \subseteq \mathbb{R} \subseteq \mathbb{C}
\]
are provided as axioms \texttt{nn-subset-zzplus},
\texttt{zzplus-subset-zz}, \texttt{zz-subset-qq},
\texttt{qq-subset-rr}, \texttt{rr-subset-cc}.  The set
\texttt{ZZplus} is also closed under \texttt{+}, \texttt{*}, and
\texttt{succ} (axioms \texttt{zzplus-add-closed},
\texttt{zzplus-mul-closed}, \texttt{zzplus-succ-closed}).

\subsubsection{Arithmetic operations}

Arithmetic operations are shared across all number systems: there is one
\texttt{+}, one \texttt{*}, one \texttt{-}, and so on.  Membership
determines which system an argument belongs to; closure axioms guarantee
the result stays in the appropriate system.  For example,
\texttt{forall([x in NN, y in NN], x + y in NN)} is the
closure axiom for \texttt{+} on \texttt{NN}.

\begin{center}
\begin{tabular}{lll}
\toprule
System & Operations & Distinguished constants \\
\midrule
\texttt{NN} &
  \texttt{+}, \texttt{*}, \texttt{succ}, \texttt{<=} &
  \texttt{1} \\[3pt]
\texttt{ZZ} &
  \texttt{+}, \texttt{*}, \texttt{-} (unary), \texttt{<=} &
  \texttt{0}, \texttt{1} \\[3pt]
\texttt{QQ} &
  \texttt{+}, \texttt{*}, \texttt{-}, \texttt{recip}, \texttt{<=} &
  \texttt{0}, \texttt{1} \\[3pt]
\texttt{RR} &
  \texttt{+}, \texttt{*}, \texttt{-}, \texttt{recip}, \texttt{abs}, \texttt{<=} &
  \texttt{0}, \texttt{1} \\[3pt]
\texttt{CC} &
  \texttt{+}, \texttt{*}, \texttt{-}, \texttt{recip}, \texttt{conjugate} &
  \texttt{0}, \texttt{1} \\
\bottomrule
\end{tabular}
\end{center}

\begin{remark}\label{remark:27135-34942-842417}
%
The infix notation \texttt{x \^{} n} is the preferred way to write
exponentiation; it parses to the same internal form as
\texttt{power(x, n)}.  Since \texttt{FUN}$(A,B)$ in VNB always means
\emph{total} on $A$, the only typing axiom is
\[
  \texttt{power} \in \mathit{FUN}(\mathit{CC} \times \mathit{NN},\, \mathit{CC})
  \qquad\text{(axiom \texttt{power-typing-nonneg})}
\]
The defining axioms are \texttt{power-zero} ($x^0 = 1$) and
\texttt{power-succ} ($x^{\mathrm{succ}(n)} = x \cdot x^n$ for
$n \in \mathit{NN}$).  A conditional equation \texttt{power-neg}
covers negative integer exponents: $x^{-n} = \mathrm{recip}(x^n)$
for $n \in \mathit{NN}$, $x \neq 0$; but this does not extend the
\texttt{FUN} typing to $\mathit{ZZ}$, since $0^{-1}$ is undefined.
Complex exponentiation $x^z$ for $z \in \mathit{CC}$ will be defined
later via the power series for \texttt{exp}.
%
\end{remark}
%
\begin{remark}\label{remark:27135-34995-146430}
%
In particular, $0^0 = 1$.  This is the mathematically correct value
for combinatorial and power-series purposes, as Knuth argues, and it is
what MIT~Scheme's exact arithmetic returns.  GNU/Maxima raises an error
here.  IMPS treated $0^0$ as undefined, which is logically sound but
forces theorems such as the binomial theorem to carry an explicit
$x \neq 0$ or $y \neq 0$ hypothesis to exclude the special case;
this system avoids that overhead by defining $0^0 = 1$.
%
\end{remark}

The arithmetic decision procedure \texttt{vnb-do-arith} (and the
\texttt{(arith)} command) evaluates \texttt{x \^{} n} on ground
numeric arguments.

\noindent
Integer literals are constants directly: \texttt{0}, \texttt{1},
\texttt{7}, \texttt{-3}, etc.\ are elements of \texttt{ZZ} (and of
\texttt{QQ}, \texttt{RR}, \texttt{CC} by inclusion); positive integer
literals are additionally elements of \texttt{NN}.  No subscripted
forms (\texttt{0\_ZZ}, \texttt{1\_RR}, etc.) are used.  In MIT~Scheme
numeric literals are first-class objects, so this representation is
natural.

Integer literals must be in canonical decimal form: no leading zeros
(other than \texttt{0} itself), so \texttt{007} and \texttt{000075} are
rejected by the parser.  \texttt{-0} normalises to \texttt{0}.

The ordinal operators \texttt{<=\_ORD}, \texttt{<\_ORD}, and
\texttt{succ\_ORD} retain their subscripts to distinguish them from the
numeric \texttt{<=} and \texttt{succ}.

\subsubsection{Axioms installed}

Each number system comes with: a sethood axiom, inclusion from the
previous system, closure of each operation, and algebraic identities
appropriate to its structure.

\begin{description}
\item[\texttt{NN}] Commutative semiring with unit \texttt{1}; no additive identity
\item[\texttt{ZZ}] Commutative ring; additive group with \texttt{0} and unary \texttt{-}
\item[\texttt{QQ}] Ordered field; \texttt{recip} defined for nonzero elements
\item[\texttt{RR}] Ordered field with \texttt{abs}; triangle inequality; \texttt{abs} zero iff argument is zero
\item[\texttt{CC}] Field; \texttt{conjugate} is an involution compatible with $+$ and $\cdot$;
  $a \cdot \overline{a} \in \mathbb{R}$ and $a \cdot \overline{a} \geq 0$
\end{description}

Because there is a single \texttt{+}, \texttt{*}, etc.\ shared across all
systems, no compatibility axioms between adjacent systems are needed:
there is nothing to be compatible with.

\begin{remark}
Each of \texttt{NN}, \texttt{ZZ}, and \texttt{QQ} carries an infinite
family of membership axioms, one per numeral appearing in the theory.
When a positive-integer numeral $n$ is introduced, it is implicitly
axiomatized as a iterated sum $1 + 1 + \cdots + 1$ ($n$ times), so
$n \in \mathtt{NN}$.  For \texttt{ZZ}, negative numerals are similarly
axiomatized as the additive inverse of such sums, giving $-n \in
\mathtt{ZZ}$.  These axioms are generated on demand: entering a numeral
literal in a formula automatically extends the theory with the
corresponding membership fact.
\end{remark}

\subsubsection{Example use in a proof}

\begin{VNBcode}
; Prove: 0 is in RR
(sp (make-wff-from-string "0 in RR"))
(ta 'rr-zero-in)                ; introduce theorem (0 in rr)
(ass)                           ; goal is now in assumptions

; Prove: for reals a, b, (a + b) is in RR
(sp (make-wff-from-string "a in RR and b in RR implies a + b in RR"))
(di)                              ; assume (a in rr) and (b in rr)
(ai "a in RR and b in RR")       ; split into (a in rr) and (b in rr)
(ta 'rr-add-closed)               ; closure axiom: forall([a, b], ...)
(inst "forall([a, b], a in rr and b in rr implies a + b in rr)" 'a)
(inst "forall([b], a in rr and b in rr implies a + b in rr)" 'b)
(bc "a in rr and b in rr implies a + b in rr")
(ass)
; => Proof complete.
\end{VNBcode}

%% =========================================================
\subsection{Ordinals}
\label{sec:ordinals}

\subsubsection{The class ORD}

\texttt{ORD} is a primitive class constant, declared in the same style as
\texttt{NN}, \texttt{RR}, and \texttt{CC}: no explicit construction of
ordinals is needed.  The class is characterised entirely by its axioms.

\begin{description}
\item[\texttt{ORD}] the class of all ordinals (a proper class)
\item[\texttt{$\alpha$ in ORD}] $\alpha$ is an ordinal
\end{description}

\subsubsection{The ordering predicate}

Ordinals are equipped with a primitive ordering predicate \texttt{<=\_ORD}
(and its strict companion \texttt{<\_ORD}) that is independent of set
membership.  The ordering has no required connection to $\in$.

\begin{description}
\item[\texttt{(<=\_ORD $\alpha$\ $\beta$)}] $\alpha \leq \beta$ in the ordinal ordering
\item[\texttt{(<\_ORD $\alpha$\ $\beta$)}] $\alpha < \beta$ (strict)
\item[\texttt{(succ\_ORD $\alpha$)}] successor ordinal of $\alpha$
\end{description}

\subsubsection{Axioms}

The following axioms are installed in \texttt{ordinals.scm}.

\begin{description}
\item[\texttt{burali-forti}]
  $\mathit{ORD} \notin \mathit{SET}$: \texttt{ORD} is a proper class.
  Every \emph{set} of ordinals has a supremum in \texttt{ORD}
  (axiom \texttt{sup-ord-in}); the supremum need not lie in the set itself
  (e.g.\ $\sup \mathit{NN} = \omega \notin \mathit{NN}$).
  If \texttt{ORD} were a set, its supremum $\alpha$ would be an ordinal,
  hence $\mathit{succ\_ORD}(\alpha) \in \mathit{ORD}$ would strictly
  exceed $\alpha$---contradicting that $\alpha$ bounds \texttt{ORD}.
\item[\texttt{ord-le-closure}]
  $\alpha \leq_O \beta \;\Rightarrow\; \alpha \in \mathit{ORD} \land \beta \in \mathit{ORD}$
\item[\texttt{ord-le-refl}]
  $\alpha \in \mathit{ORD} \;\Rightarrow\; \alpha \leq_O \alpha$
\item[\texttt{ord-le-antisymm}]
  $\alpha \leq_O \beta \land \beta \leq_O \alpha \;\Rightarrow\; \alpha = \beta$
\item[\texttt{ord-le-trans}]
  $\alpha \leq_O \beta \land \beta \leq_O \gamma \;\Rightarrow\; \alpha \leq_O \gamma$
\item[\texttt{ord-le-total}]
  $\alpha,\beta \in \mathit{ORD} \;\Rightarrow\; \alpha \leq_O \beta \lor \beta \leq_O \alpha$
\item[\texttt{ord-zero-least}]
  $\alpha \in \mathit{ORD} \;\Rightarrow\; 0 \leq_O \alpha$
\item[\texttt{ord-lt-iff}]
  $\alpha <_O \beta \;\Leftrightarrow\; \alpha \leq_O \beta \land \alpha \neq \beta$
\item[\texttt{ord-succ-in}]
  $\alpha \in \mathit{ORD} \;\Rightarrow\; \mathit{succ\_ORD}(\alpha) \in \mathit{ORD}$
\item[\texttt{ord-succ-above}]
  $\alpha \in \mathit{ORD} \;\Rightarrow\; \alpha <_O \mathit{succ\_ORD}(\alpha)$
\item[\texttt{ord-succ-immediate}]
  $\alpha,\beta \in \mathit{ORD},\; \alpha <_O \beta \;\Rightarrow\; \mathit{succ\_ORD}(\alpha) \leq_O \beta$
\item[\texttt{ord-succ-injective}]
  $\mathit{succ\_ORD}(\alpha) = \mathit{succ\_ORD}(\beta) \;\Rightarrow\; \alpha = \beta$
\item[\texttt{limit-ord-iff}]
  $\mathit{LIMIT\text{-}ORD}(\lambda) \;\Leftrightarrow\; \lambda \in \mathit{ORD} \land \lambda \neq 0 \land \lnot\exists\alpha.\,\lambda = \mathit{succ\_ORD}(\alpha)$
\item[\texttt{ord-segment-is-set}]
  $\alpha \in \mathit{ORD} \;\Rightarrow\; \mathit{ORD\text{-}SEGMENT}(\alpha) \in \mathit{SET}$
\item[\texttt{ord-segment-membership}]
  $x \in \mathit{ORD\text{-}SEGMENT}(\alpha) \;\Leftrightarrow\; x <_O \alpha$
\item[\texttt{ord-segment-zero}]
  $\mathit{ORD\text{-}SEGMENT}(0) = \emptyset$
\item[\texttt{ord-segment-succ}]
  $x \in \mathit{ORD\text{-}SEGMENT}(\mathit{succ\_ORD}(\alpha)) \;\Leftrightarrow\; x \in \mathit{ORD\text{-}SEGMENT}(\alpha) \lor x = \alpha$
\item[\texttt{sup-ord-in}]
  $A \in \mathit{SET},\; A \subseteq \mathit{ORD} \;\Rightarrow\; \mathit{SUP\text{-}ORD}(A) \in \mathit{ORD}$
\item[\texttt{sup-ord-upper}]
  $\mathit{SUP\text{-}ORD}(A)$ is an upper bound of $A$
\item[\texttt{sup-ord-least}]
  $\mathit{SUP\text{-}ORD}(A)$ is the least upper bound
\item[\texttt{sup-ord-empty}]
  $\mathit{SUP\text{-}ORD}(\emptyset) = 0$
\item[\texttt{transfinite-induction}]
  see \S\ref{sec:ord-induction}
\end{description}

Here $\leq_O$ and $<_O$ abbreviate \texttt{<=\_ORD} and \texttt{<\_ORD}.

\subsubsection{Relationship to number systems}

\texttt{NN} is a subset of \texttt{ORD} (axiom \texttt{nn-subset-ord}):
\[
  \forall n.\; n \in \mathit{NN} \;\Rightarrow\; n \in \mathit{ORD}
\]
Consequently $0 \in \mathit{ORD}$ (from \texttt{nn-zero-in} and
\texttt{nn-subset-ord}).  The ordinal successor agrees with the numeric
successor on $\mathit{NN}$ (axiom \texttt{ord-succ-nn}), and
\texttt{<=\_ORD} restricted to $\mathit{NN}$ coincides with the numeric
\texttt{<=} (axiom \texttt{ord-le-nn-compat}).  The number systems
\texttt{ZZ}, \texttt{QQ}, \texttt{RR}, \texttt{CC} are not subsets of
\texttt{ORD}.

\subsubsection{Transfinite induction}
\label{sec:ord-induction}

\paragraph{The induction axiom.}

The axiom \texttt{transfinite-induction} is the \emph{complete} (strong)
form: a class $C$ that contains every ordinal whose strict predecessors are
all in $C$ must contain every ordinal.
\begin{align*}
  &\forall C.\;
    \bigl(\forall\alpha \in \mathit{ORD}.\;
      (\forall \beta <_O \alpha.\; \beta \in C) \;\Rightarrow\; \alpha \in C\bigr) \\
  &\quad\Rightarrow\;
    \forall\alpha \in \mathit{ORD}.\; \alpha \in C
\end{align*}

\paragraph{The \texttt{(tfi)} proof command.}

\texttt{(tfi)} applies transfinite induction to the current focus goal.
The goal must have the form
\[
  \forall \alpha.\; \alpha \in \mathit{ORD} \;\Rightarrow\; P(\alpha)
\]
and \texttt{(tfi)} replaces it with the single inductive-step subgoal
\[
  \forall \alpha.\;
    \bigl(\alpha \in \mathit{ORD} \;\land\;
          \forall\beta <_O \alpha.\; P(\beta)\bigr)
    \;\Rightarrow\; P(\alpha)
\]

\paragraph{The \texttt{(tfi3)} proof command.}

\texttt{(tfi3)} splits the same goal into three subgoals matching the
standard Peano-style cases:
\begin{enumerate}
\item \textbf{Base:} $P(0)$.
\item \textbf{Successor:} $\forall\alpha \in \mathit{ORD}.\; P(\alpha) \Rightarrow P(\mathit{succ\_ORD}(\alpha))$.
\item \textbf{Limit:} $\forall\lambda.\;
  \mathit{LIMIT\text{-}ORD}(\lambda) \land
  (\forall\beta <_O \lambda.\; P(\beta)) \;\Rightarrow\; P(\lambda)$.
\end{enumerate}

\subsubsection{Transfinite recursion}
\label{sec:ord-recursion}

The theoretical basis for transfinite recursive definitions is the
following pair of propositions.  Both are proved by the standard
least-counterexample argument using well-foundedness of $\leq_O$.

\begin{prop}[Transfinite recursion, set version]\label{prop:tfrec-set}
Let $U = \mathit{ORD\text{-}SEGMENT}(\alpha)$ for some $\alpha \in \mathit{ORD}$,
let $A \in \mathit{SET}$, and let
\[
  \mathcal{F}(U, A) \;=\; \bigcup_{x \in U}\,\mathit{FUN}(\mathit{ORD\text{-}SEGMENT}(x),\, A)
\]
be the set of all functions defined on an initial segment within $U$ with
values in $A$.  For any $\mathcal{E} \in \mathit{FUN}(\mathcal{F}(U,A),\, A)$ there
is a unique $f \in \mathit{FUN}(U, A)$ such that
\[
  f(x) \;=\; \mathcal{E}\bigl(f \!\restriction\! \mathit{ORD\text{-}SEGMENT}(x)\bigr)
  \quad \text{for every } x \in U.
\]
\end{prop}

\begin{prop}[Transfinite recursion, class version]\label{prop:tfrec-class}
Let $A$ be a class, and let $\mathcal{E}$ be a function from the class of
all functions defined on initial ordinal segments with values in $A$ to $A$.
Then there is a unique function $f$ on $\mathit{ORD}$ with values in $A$ such that
\[
  f(\xi) \;=\; \mathcal{E}\bigl(f \!\restriction\! \mathit{ORD\text{-}SEGMENT}(\xi)\bigr)
  \quad \text{for every } \xi \in \mathit{ORD}.
\]
\end{prop}

\begin{remark}
In Proposition~\ref{prop:tfrec-class} the function $f$ has $\mathit{ORD}$ as its
domain, which is a proper class.  Such $f$ is not an element of $\mathit{FUN}(A,B)$
(which requires set-valued domains), but it is still a well-defined class of
ordered pairs with the function property.
\end{remark}

In practice the step function $\mathcal{E}$ is specified by splitting on
three mutually exclusive cases for the argument $\xi$:
\begin{itemize}[label=$\circ$, leftmargin=0.9cm]
\item \textbf{Base:} $\xi = 0$
\item \textbf{Successor:} $\xi = \mathit{succ\_ORD}(\alpha)$ for a unique predecessor $\alpha$
\item \textbf{Limit:} $\mathit{LIMIT\text{-}ORD}(\xi)$, i.e.\ $\xi \neq 0$ and $\xi$ is not a successor
\end{itemize}

The Scheme helper \texttt{def-by-ord-recursion} implements exactly this
pattern (Proposition~\ref{prop:tfrec-class}): it defines a new function
constant $F$ by the three standard cases and installs the defining
equations as named axioms.

\begin{VNBcode}
(def-by-ord-recursion 'F
  base-val            ; F(0) = base-val
  '(alpha val)        ; successor: F(succ_ORD(alpha)) = succ-expr
  succ-expr           ;   where val stands for F(alpha)
  '(lambda)           ; limit: F(lambda) = lim-expr
  lim-expr)           ;   where lambda is the limit ordinal
\end{VNBcode}

Three axioms are installed:
\begin{description}
\item[\texttt{F-zero}] \texttt{F(0) = base-val}
\item[\texttt{F-succ}] \texttt{forall([alpha in ORD], F(succ\_ORD(alpha)) = succ-expr[val:=F(alpha)])}
\item[\texttt{F-limit}] \texttt{forall([lambda], LIMIT-ORD(lambda) implies F(lambda) = lim-expr)}
\end{description}

\paragraph{Example: identity function on ordinals.}

\begin{VNBcode}
(def-by-ord-recursion 'COPY-ORD
  0                                         ; COPY-ORD(0) = 0
  '(alpha val) '(succ_ORD val)              ; COPY-ORD(succ(a)) = succ(COPY-ORD(a))
  '(lam) '(SUP-ORD (ORD-SEGMENT lam)))      ; COPY-ORD(lam) = sup of predecessors
\end{VNBcode}

\paragraph{Example: ordinal addition $\alpha + \beta$ (fixing $\alpha$).}

Defining $A_\alpha(\beta) = \alpha + \beta$ as a function of $\beta$:
\begin{VNBcode}
(def-by-ord-recursion 'ORD-ADD
  'alpha-fixed                              ; ORD-ADD(0) = alpha-fixed
  '(beta val) '(succ_ORD val)              ; ORD-ADD(succ(b)) = succ(ORD-ADD(b))
  '(lam) '(SUP-ORD (ORD-SEGMENT lam)))     ; ORD-ADD(lam) = sup of ORD-ADD on lam
\end{VNBcode}

(Here \texttt{alpha-fixed} is whatever constant plays the role of
$\alpha$.  A full definition of addition would use a nested recursion.)

%% =========================================================
\section{Theories and axioms}
\label{sec:axioms}

The base theory \texttt{vnb-set-theory} is loaded automatically and
provides the following named axioms (all usable via
\texttt{cmd-theorem-assumption}):

\begin{description}
\item[\texttt{membership-implies-sethood}]
  $\forall a\,\forall b.\;(a \in b) \Rightarrow a \in \mathit{SET}$
\item[\texttt{extensionality}]
  $\forall a\,\forall b.\;(a \in \mathit{SET} \land b \in \mathit{SET}) \Rightarrow
  (a = b \Leftrightarrow \forall x.\,(x \in a \Leftrightarrow x \in b))$
\item[\texttt{empty-set-is-set}]
  $\texttt{EMPTY-SET} \in \mathit{SET}$
\item[\texttt{empty-set-has-no-members}]
  $\forall x.\;\lnot(x \in \texttt{EMPTY-SET})$
\item[\texttt{pairing}]
  $\forall a\,\forall b.\;\{a,b\} \in \mathit{SET}$
\item[\texttt{pairing-membership}]
  $\forall a\,\forall b\,\forall x.\;(x \in \{a,b\}) \Leftrightarrow (x=a \lor x=b)$
\item[\texttt{power-set}]
  $\forall a.\;a \in \mathit{SET} \Rightarrow \texttt{power}(a) \in \mathit{SET}$
\item[\texttt{power-set-membership}]
  $\forall a\,\forall x.\;(x \in \texttt{power}(a)) \Leftrightarrow
  (x \in \mathit{SET} \land \forall z.\,(z \in x \Rightarrow z \in a))$
\end{description}

To add a new theorem to the theory (and automatically install its macete):
\begin{VNBcode}
(theory-add-theorem! (current-theory) 'my-thm
  (make-wff-from-string "forall([x in A], x in B)"))
\end{VNBcode}

\subsection{Defined constants}
\label{sec:def-constant}

A \emph{defined constant} is a new constant symbol introduced by a
collection of characterizing axioms.  It differs from a plain
\texttt{theory-add-axiom!} call only in bookkeeping: the axioms are
recorded under a separate \texttt{definitions} slot in the theory record,
so the system can distinguish what was \emph{postulated} from what was
\emph{introduced by definition}.

\begin{VNBcode}
(def-constant 'NAME
  '(axiom-name-1 formula-1)
  '(axiom-name-2 formula-2)
  ...)
\end{VNBcode}

Each extra argument is a two-element list \texttt{(axiom-name formula)}.
The formulas are installed in the global theorem and macete tables exactly
as \texttt{theory-add-axiom!} would install them, so they are immediately
available to proof commands and macetes.

\textbf{Existence and uniqueness are the user's responsibility.}  The
system accepts the definition without checking that the described object
exists or is unique.  For constants where this matters (e.g., a least
element), you should prove existence and uniqueness before invoking
\texttt{def-constant}, or at least note the obligation explicitly.

Example---introducing the first uncountable ordinal $\omega_1$:
\begin{VNBcode}
(def-constant 'OMEGA1
  '(omega1-is-ordinal
    (IN OMEGA1 ORD))
  '(omega1-is-uncountable
    (NOT (COUNTABLE OMEGA1)))
  '(omega1-is-least-unc
    (FORALL alpha
      (IMPLIES (AND (IN alpha ORD) (NOT (COUNTABLE alpha)))
               (<= OMEGA1 alpha)))))
\end{VNBcode}

To display all defined constants and their characterizing axioms:
\begin{VNBcode}
(display-definitions)
\end{VNBcode}

\subsection{Pending axioms}
\label{sec:pending-axioms}

The following axioms are needed for full set-theoretic reasoning but are
not yet installed in the base theory.

\begin{description}

\item[\texttt{list-sethood}]
  $[e_1,\ldots,e_n] \in \mathit{SET}$ for all $e_i$;
  needed whenever a tuple must itself be proven a set.

\item[\texttt{separation}]
  $A \in \mathit{SET} \Rightarrow \{x \in A \mid p\} \in \mathit{SET}$;
  required to use a separation class as a membership target, and to prove
  that the intersection of a class with a set is a set.

\item[\texttt{union-set}]
  $A \in \mathit{SET} \Rightarrow \texttt{UNION-SET}(A) \in \mathit{SET}$
  (big union: $\bigcup A = \{x \mid \exists a \in A.\; x \in a\}$).
  Distinct from the n-ary \texttt{union} constructor.

\item[\texttt{union-set-membership}]
  $x \in \texttt{UNION-SET}(A) \Leftrightarrow \exists a.\;(a \in A \land x \in a)$.

\item[\texttt{cartesian-n-ary-sethood}]
  $n$-ary version: all $A_i \in \mathit{SET} \Rightarrow
  A_1 \times \cdots \times A_n \in \mathit{SET}$.  The binary case is
  installed (\texttt{cartesian-set-iff}); $n$-ary cases reduce to the
  binary one by associativity but require explicit axioms or a derivation.

\item[\texttt{tuples-induction}]
  Structural induction over $\mathit{TUPLES}(A)$: to prove $P(s)$ for all
  $s \in \mathit{TUPLES}(A)$, it suffices to prove $P([\,])$ and to prove
  $P([a, a_1, \ldots, a_n])$ assuming $a \in A$ and $P([a_1,\ldots,a_n])$.
  The membership rule for the intro/elim commands is installed as a
  primitive inference; this induction schema is pending.

\item[\texttt{prepend} and \texttt{length-recursive}]
  A \texttt{prepend}($x$, $L$) constructor (prepend element $x$ to tuple
  $L$) is needed to state the recursive characterisation of length:
  $\texttt{length}(\texttt{prepend}(x, L)) = \texttt{length}(L) + 1$.
  The base axioms (\texttt{length-of-empty}, \texttt{length-in-nn},
  \texttt{nth-in-range}) are already installed; the recursive step requires
  \texttt{prepend} and \texttt{tuples-induction}.

\item[\texttt{fun-domain-apply-def}]
  $(f \in \mathit{FUN}(A)) \land (x \in A) \Rightarrow f(x) = f(x)$
  (definedness of application).

\item[\texttt{fun-domain-extensionality}]
  $f,g \in \mathit{FUN}(A) \land (\forall x \in A.\;f(x)=g(x))
  \Rightarrow f = g$.

\item[\texttt{fun-codomain-iff}]
  $f \in \mathit{FUN}(A,B) \;\Leftrightarrow\;
  (f \in \mathit{FUN}(A)) \land (\forall x \in A.\;f(x)\in B)$.

\item[\texttt{fun-apply-type}]
  $(f \in \mathit{FUN}(A,B)) \land (x \in A) \Rightarrow f(x) \in B$
  (derivable from \texttt{fun-codomain-iff}).

\item[\texttt{equality-symmetry}]
  $a = b \Rightarrow b = a$.

\item[\texttt{equality-transitivity}]
  $a = b \land b = c \Rightarrow a = c$.

\item[\texttt{equality-substitution}]
  $a = b \land P(a) \Rightarrow P(b)$ (for any formula $P$).

\item[\texttt{quasi-eq-reflexivity}]
  $(a == a)$ for all $a$ (logical axiom).

\item[\texttt{quasi-eq-symmetry}]
  $(a == b) \Rightarrow (b == a)$.

\item[\texttt{quasi-eq-transitivity}]
  $(a == b) \land (b == c) \Rightarrow (a == c)$.

\item[\texttt{quasi-eq-def}]
  $(a = a) \Rightarrow ((a == b) \Leftrightarrow (a = b))$;
  when $a$ is defined, quasi-equality coincides with strict equality.

\item[\texttt{nn-subset-ord}]
  $\forall n.\; n \in \mathit{NN} \Rightarrow n \in \mathit{ORD}$.

\item[\texttt{ord-def} and related]
  Ordering, induction, successor, and properness (\texttt{burali-forti})
  axioms for \texttt{ORD}; see Section~\ref{sec:ordinals}.

\item[\texttt{infinity}]
  There exists an infinite set.

\item[\texttt{choice}]
  $A \neq \emptyset \Rightarrow \texttt{choice}(A) \in A$.
  Global primitive choice operator; see
  Section~\ref{sec:choice-iota-lambda}.

\item[\texttt{comp-membership}]
  $x \in \{y \mid p\} \Leftrightarrow p[y \mapsto x]$;
  class comprehension membership.

\item[\texttt{vnb-lambda-type}]
  $(\forall x \in A.\;\mathit{body}(x) \in B) \Rightarrow
   \texttt{vnb-lambda}(x, \mathit{body}) \in \mathit{FUN}(A,B)$.

\item[\texttt{vnb-lambda-beta}]
  $\texttt{vnb-lambda}(x, \mathit{body})(v) = \mathit{body}[x \mapsto v]$.

\item[\texttt{iota-def}]
  $(\exists! x.\;p(x)) \Rightarrow p[\texttt{iota}(x, p)]$
  and the result is the unique such element.

\end{description}

\noindent
Currently only reflexivity ($a = a$) is available as a primitive inference
(\texttt{cmd-reflexivity}).  Without the equality rules above, proofs of
equalities other than reflexive ones must be carried out indirectly.

%% =========================================================
\section{Mathematical structure definitions}

\subsection{The \texttt{def-structure} form}

\texttt{def-structure} declares a named mathematical structure at the
Scheme level.  It registers the structure, installs NTH-based accessor
macetes for every carrier set and every operation, and posts a
definitional axiom for the associated \texttt{IS-\textit{NAME}} predicate
in the current theory.

\begin{VNBcode}
(def-structure name carriers op-specs axiom-names)
\end{VNBcode}

\begin{description}

\item[\texttt{name}]
  Quoted symbol naming the structure, e.g.\ \texttt{'BANACHSPACE}.

\item[\texttt{carriers}]
  Quoted list of accessor name symbols for the carrier sets,
  e.g.\ \texttt{'(VECTORS)}.  A structure with $m$ carriers
  $C_1, \ldots, C_m$ has its $i$-th carrier at \texttt{nth($i$, s)} in
  any instance~\texttt{s}.

\item[\texttt{op-specs}]
  Quoted list of operation and constant specifications.  Two forms are
  accepted:
  \begin{itemize}\setlength\itemsep{2pt}
  \item \emph{Function spec} \texttt{($\mathit{op}$\ $\mathit{domain}$\ $\mathit{range}$)}
    (three elements): asserts
    $\mathit{op}(s) \in \mathit{FUN}(\mathit{domain},\mathit{range})$.
  \item \emph{Constant spec} \texttt{($\mathit{op}$\ $\mathit{set}$)}
    (two elements): asserts $\mathit{op}(s) \in \mathit{set}$ (an
    element, not a function).
  \end{itemize}
  Carrier and op accessor names appear \emph{bare}; \texttt{def-structure}
  expands them to \texttt{($\mathit{accessor}$\ s)}.  External constants
  such as \texttt{CC} and \texttt{RR} are left unchanged.  All ops are
  indexed after the carriers in NTH order.

\item[\texttt{axiom-names}]
  Quoted list of theorem names that axiomatize the structure
  (may be \texttt{'()}).

\end{description}

\subsection{What \texttt{def-structure} installs}

\paragraph{Accessor macetes.}
For each carrier $C_i$ (at NTH index $i$) and each operation
$\mathit{op}_j$ (at NTH index $m+j$), a macete is installed that
rewrites the accessor term to the corresponding NTH projection:
\[
  (C_i\;\; s) \;\longmapsto\; (\texttt{NTH}\;\; i\;\; s)
  \qquad
  (\mathit{op}_j\;\; s) \;\longmapsto\; (\texttt{NTH}\;\; m{+}j\;\; s)
\]
The schema variable matches any expression, so the macete fires anywhere
in a goal.

\paragraph{IS-\textit{NAME} axiom.}
The predicate \texttt{IS-\textit{NAME}} is axiomatized as a conjunction
of typing conditions on each accessor:
\begin{align*}
  \forall s.\;\mathit{IS\text{-}NAME}(s) \;\Leftrightarrow\; &
  C_1(s) \in \mathit{SET} \land \cdots \land C_m(s) \in \mathit{SET}\\
  &\land\; \mathit{op}_1(s) \in \mathit{FUN}(\mathit{dom}_1, \mathit{rng}_1)
  \land \cdots
\end{align*}
The axiom is named \texttt{IS-\textit{NAME}} and is immediately available
via \texttt{cmd-theorem-assumption} and \texttt{cmd-apply-macete}.

\subsection{Example: Banach spaces over \texorpdfstring{$\mathbb{C}$}{CC}}

\begin{VNBcode}
(def-structure 'BANACHSPACE
  '(VECTORS)
  '((SCALAR-TIMES (CARTESIAN CC VECTORS) VECTORS)
    (PLUS         (CARTESIAN VECTORS VECTORS) VECTORS)
    (NORM         VECTORS RR))
  '())
\end{VNBcode}

This installs four accessor macetes (VECTORS at index~1, SCALAR-TIMES at
2, PLUS at 3, NORM at 4) and adds the following axiom to the current
theory (here \texttt{s} is the FORALL-bound instance variable in the
generated formula, not something the user writes):

\begin{VNBcode}
(FORALL s
  (IFF (IS-BANACHSPACE s)
       (AND (IN (VECTORS s) SET)
            (AND (IN (SCALAR-TIMES s)
                     (FUN (CARTESIAN CC (VECTORS s)) (VECTORS s)))
                 (AND (IN (PLUS s)
                          (FUN (CARTESIAN (VECTORS s) (VECTORS s))
                               (VECTORS s)))
                      (IN (NORM s) (FUN (VECTORS s) RR)))))))
\end{VNBcode}

Working with a Banach space hypothesis in a proof:
\begin{VNBcode}
; Assume (IS-BANACHSPACE E); unfold its definition
(cmd-apply-macete ps 'IS-BANACHSPACE)
; Goal rewritten to: (AND (IN (VECTORS E) SET) ...)

; Reduce any accessor term to its NTH form
(cmd-apply-macete ps 'VECTORS)
; Every (VECTORS E) in the goal becomes (NTH 1 E)
\end{VNBcode}

\subsection{Further structure definitions}

The same mechanism handles any arity of carriers and operations.  For
instance, a vector space over an arbitrary field~$F$:

\begin{VNBcode}
(def-structure 'VECTORSPACE
  '(FIELD VECTORS)                 ; two carriers
  '((SCALAR-TIMES (CARTESIAN FIELD VECTORS) VECTORS) ; function spec
    (PLUS         (CARTESIAN VECTORS VECTORS) VECTORS); function spec
    (ZERO         VECTORS)                            ; constant spec: zero vector
    (NEG          VECTORS VECTORS))                   ; function spec
  '())
; FIELD    -> (NTH 1 s),  VECTORS -> (NTH 2 s)
; SCALAR-TIMES -> (NTH 3 s),  PLUS -> (NTH 4 s)
; ZERO     -> (NTH 5 s),  NEG  -> (NTH 6 s)
\end{VNBcode}

\section{\texttt{def-functoid}}
\label{sec:def-functoid}

\texttt{def-functoid} introduces a named functoid at the proof-checker
level.  It is a definitional device, not an object of the theory: the
name it installs can be applied and reasoned about, but there is no
formal quantifier ranging over functoids.

The general form is:
\begin{VNBcode}
(def-functoid name (x) body)
\end{VNBcode}
where \texttt{x} is a formal parameter ranging over classes and
\texttt{body} is a class-valued expression, typically built from
accessor macetes installed by \texttt{def-structure}.

\paragraph{Example: the forgetful functor.}
The forgetful functor from normed complex vector spaces to their
underlying vector spaces drops the norm component:

\begin{VNBcode}
(def-functoid NormedSpace2ComplexVecSpace (S)
  [(VECTORS S), (SCALAR-TIMES S), (PLUS S), (ZERO S), (NEG S)])
\end{VNBcode}

The domain of this functoid is the proper class of all normed complex
vector spaces; the codomain is the proper class of all complex vector
spaces.  Neither is a set, which is precisely why \texttt{fun} cannot
express this mapping and a functoid is required.

\paragraph{Example: the carrier functor.}
\begin{VNBcode}
(def-functoid GroupCarrier (G) (CARRIER G))
\end{VNBcode}
This maps every group to its underlying carrier set.  The result is
always a set (\texttt{IS-GROUP} entails \texttt{(IN (CARRIER G) SET)}),
even though the domain is a proper class.

%% =========================================================
\chapter{Inference mechanisms}

\section{Sequents and deduction graphs}\label{section:sequents}

A \emph{sequent} is a structure consisting of an ordered pair $(L,
B)$, where $L$ A list of formulas called \emph{assumptions} \and $B$
is a single formula called the \emph{assertion}. A sequent is
represented as follows
\[
  A_1,\, A_2,\, \ldots,\, A_n \;\Longrightarrow\; B
\]
where $A_1, \ldots, A_n$ are the assumptions (or hypotheses or the context) and $B$
is the assertion. In the course of a proof, many sequents may be created.

A \emph{deduction graph} is a directed acyclic graph of
\emph{sequent nodes} connected by \emph{inference nodes}.  An inference
node records which rule was applied, which sequent nodes are its
hypotheses, and which sequent node is its conclusion.

A sequent node is \emph{grounded} when it has an inference node all of
whose hypothesis nodes are themselves grounded.  Leaf inferences (with no
hypothesis nodes) are immediately grounded.  Grounding propagates upward.
A proof is \emph{complete} when the root sequent node is grounded.

\begin{remark}
A \emph{proof state} (\texttt{<proof-state>}) is not the same thing as a
deduction graph.  It is a thin wrapper that bundles three items together: the
deduction graph itself, the \emph{root} sequent node (the original goal,
fixed at \texttt{start-proof} time and checked by \texttt{qed}), and the
\emph{focus} sequent node (the subgoal currently being worked on).  The
deduction graph is the mathematical object that grows as rules are applied;
the proof state is the navigational view into it that the interactive session
and the \texttt{cmd-*} commands operate on.  In the interactive session,
\texttt{*ps*} always holds the current proof state.
\end{remark}

%% =========================================================
\section{Starting and inspecting a proof}

\begin{VNBcode}
; Start a new proof of a formula
(define ps (start-proof (make-wff-from-string "n in NN and m in NN implies n in NN")))

; Display the current state
(print-proof-state ps)

; Check if done
(proof-done? ps)   ; => #f or #t

; List open goals
(proof-open-goals ps)
\end{VNBcode}

The variable \texttt{ps} holds the \emph{proof state}: a deduction graph
plus the current \emph{focus} sequent node (the subgoal currently being
worked on).  Every proof command returns a new proof state with an
updated focus.

\texttt{start-proof} takes a \texttt{wff} (produced by \texttt{make-wff} or \texttt{make-wff-from-string}).
Any local contexts attached to the \texttt{wff} are unrolled into initial
sequent assumptions, so that subsequent inference rules see them as
ordinary hypotheses.  All sequent assumptions and assertions are
\texttt{wff} objects throughout the proof; the theory name propagates
automatically to every subformula produced by inference rules.

%% =========================================================
\section{Proof commands}
\label{sec:commands}

Commands are Scheme procedures of the form
\[
  \texttt{(cmd-}\mathit{name}\;\; \mathit{ps}\;\; \mathit{arg}_1 \ldots\texttt{)}
\]
Each command applies one primitive inference to the focus node, then
advances the focus to the first new ungrounded subgoal.

\section{Logical connectives}

\begin{description}

\item[\texttt{(cmd-direct-inference ps)}]
  Decompose the goal by its top-level connective:
  \begin{itemize}[label=$-$, leftmargin=1cm]
  \item \texttt{p and q} $\to$ two subgoals $p$, $q$.
  \item \texttt{p implies q} $\to$ assume $p$, prove $q$.
  \item \texttt{p iff q} $\to$ two subgoals (each direction).
  \item \texttt{not(p)} $\to$ assume $p$, prove \texttt{FALSEHOOD}.
  \item \texttt{forall([\ldots], p)} $\to$ peel the entire leading
    \texttt{forall} chain in one step, introducing a fresh variable for
    each bound variable.  If the substituted body is
    \texttt{$y$ in $A$ implies \ldots}, absorb \texttt{$y$ in $A$} into
    the assumptions and continue peeling.  Stop at the first
    non-\texttt{forall} head.
    Example: goal \texttt{forall([x in A, y], body)} yields one subgoal
    $x_0 \in A \;\Longrightarrow\; \texttt{body}[x/x_0,\,y/y_0]$.
  \end{itemize}

\item[\texttt{(cmd-antecedent-inference ps $f$)}]
  Decompose assumption $f$ by its top-level connective:
  \begin{itemize}[label=$-$, leftmargin=1cm]
  \item \texttt{p and q} $\to$ replace with $p$ and $q$ separately.
  \item \texttt{p or q} $\to$ case split: branch on $p$, branch on $q$.
  \item \texttt{not(p)} $\to$ derive \texttt{FALSEHOOD} if $p$ also in context.
  \item \texttt{forsome([\ldots], p)} $\to$ introduce a fresh witness for $x$.
  \end{itemize}

\item[\texttt{(cmd-assumption ps)}]
  Close goal immediately if it matches an assumption or is \texttt{TRUTH}.

\item[\texttt{(cmd-arith ps)}]
  Close the current goal by ground arithmetic evaluation.  The goal must
  be a closed formula (no free variables) built from numeric literals,
  $+$, $*$, $-$, \texttt{recip}, \texttt{abs}, \texttt{succ},
  \texttt{conjugate}, $\leq$, $=$,
  $\in \mathit{NN}/\mathit{ZZ}/\mathit{QQ}/\mathit{RR}/\mathit{CC}$,
  and the propositional connectives.  Exact (rational) arithmetic only;
  no axiom chain is constructed.
  Related: \texttt{(vnb-do-arith \textit{f})} returns \texttt{\#t}/\texttt{\#f}
  for use outside a proof.

\item[\texttt{(cmd-reflexivity ps)}]
  Close goal \texttt{a = a} when \texttt{a} is demonstrably defined.
  Fails if \texttt{a} may be undefined.

\item[\texttt{(cmd-quasi-reflexivity ps)}]
  Close goal \texttt{a == a} unconditionally, for any term \texttt{a}.

\item[\texttt{(cmd-cut ps $L$)}]
  Introduce lemma $L$: subgoal~1 proves $L$; subgoal~2 has $L$ as assumption.

\item[\texttt{(cmd-or-intro-left ps)}]
  To prove \texttt{p or q}, prove $p$ instead.

\item[\texttt{(cmd-or-intro-right ps)}]
  To prove \texttt{p or q}, prove $q$ instead.

\item[\texttt{(cmd-proof-by-contradiction ps)}]
  To prove $p$, assume \texttt{not(p)} and derive \texttt{FALSEHOOD}.

\item[\texttt{(cmd-instantiate ps $F$\ $t$)}]
  From \texttt{forall([$x$], $p$)} $= F$ in context, add $p[x/t]$.

\item[\texttt{(cmd-exists-witness ps $t$)}]
  To prove \texttt{forsome([$x$], $p$)}, prove $p[x/t]$.

\item[\texttt{(cmd-weaken ps $f$)}]
  Remove formula $f$ from the assumptions.

\item[\texttt{(cmd-backchain ps $F$)}]
  From \texttt{p implies goal} $= F$ in context, generate subgoal~$p$.

\item[\texttt{(cmd-theorem-assumption ps $\mathit{name}$)}]
  Add the named theorem (or axiom) to the assumptions.

\item[\texttt{(cmd-apply-macete ps $\mathit{name}$)}]
  Apply named macete: rewrite the goal using the associated theorem.

\item[\texttt{(cmd-qed ps $\mathit{name}$)}]
  After a proof completes, install its root assertion as a named theorem
  in the current theory.  Errors if the proof is not complete.

\end{description}

\subsection{N-ary constructors}

\begin{description}

\item[\texttt{(cmd-cartesian-intro ps)}]
  To prove \texttt{[$a_1$,\ldots,$a_n$] in cartesian($A_1$,\ldots,$A_n$)},
  generate $n$ subgoals \texttt{$a_i$ in $A_i$}.

\item[\texttt{(cmd-cartesian-elim ps $F$\ $k$)}]
  From \texttt{[$a_1$,\ldots,$a_n$] in cartesian($A_1$,\ldots,$A_n$)} $= F$
  in context, add \texttt{$a_k$ in $A_k$}.

\item[\texttt{(cmd-tuples-intro ps)}]
  To prove \texttt{[$a_1$,\ldots,$a_n$] in tuples($A$)},
  generate $n$ subgoals \texttt{$a_i$ in $A$}.
  For the empty sequence \texttt{[]} the goal closes immediately (0 subgoals).

\item[\texttt{(cmd-tuples-elim ps $F$\ $k$)}]
  From \texttt{[$a_1$,\ldots,$a_n$] in tuples($A$)} $= F$
  in context, add \texttt{$a_k$ in $A$}.

\item[\texttt{(cmd-nth-reduce ps)}]
  Rewrite every occurrence of \texttt{nth($k$, [$\ldots$])} in the goal to
  the $k$-th element (1-indexed).

\item[\texttt{(cmd-vnb-lambda-beta ps)}]
  Rewrite every occurrence of \texttt{vnb-lambda($x$, $\mathit{body}$)($v$)}
  in the goal to $\mathit{body}[x := v]$ (capture-avoiding).

\item[\texttt{(cmd-union-intro ps $k$)}]
  To prove \texttt{$x$ in union($A_1$,\ldots,$A_n$)},
  prove \texttt{$x$ in $A_k$} instead.

\item[\texttt{(cmd-union-elim ps $F$)}]
  From \texttt{$x$ in union($A_1$,\ldots,$A_n$)} $= F$ in context,
  case-split: one branch per $A_i$ with \texttt{$x$ in $A_i$} assumed.

\item[\texttt{(cmd-intersection-intro ps)}]
  To prove \texttt{$x$ in intersection($A_1$,\ldots,$A_n$)},
  generate $n$ subgoals \texttt{$x$ in $A_i$}.

\item[\texttt{(cmd-intersection-elim ps $F$\ $k$)}]
  From \texttt{$x$ in intersection($A_1$,\ldots,$A_n$)} $= F$ in context,
  add \texttt{$x$ in $A_k$}.

\end{description}

%% =========================================================
\section{Interactive short-form commands}
\label{sec:interactive}

The file \texttt{interactive.scm} provides single- or two-letter wrappers
around every \texttt{cmd-*} proof command, together with the global proof
state variable \texttt{*ps*} and the \texttt{show} display function.
These are the primary commands used in a live proof session (including
through the Emacs interface).

\begin{description}

\item[\texttt{(sp \textit{wff})}]
  \texttt{start-proof} --- Start a new proof from a \texttt{wff};
  context unrolled to assumptions.

\item[\texttt{(di)}]
  \texttt{cmd-direct-inference} --- Decompose goal; peels leading
  \texttt{FORALL} chain.

\item[\texttt{(pbc)}] \texttt{cmd-proof-by-contradiction}

\item[\texttt{(ass)}] \texttt{cmd-assumption}

\item[\texttt{(arith)}]
  \texttt{cmd-arith} --- Close goal by exact arithmetic evaluation.

\item[\texttt{(rfl)}] \texttt{cmd-reflexivity}

\item[\texttt{(qrfl)}] \texttt{cmd-quasi-reflexivity}

\item[\texttt{(oi-l)}] \texttt{cmd-or-intro-left}

\item[\texttt{(oi-r)}] \texttt{cmd-or-intro-right}

\item[\texttt{(ci)}] \texttt{cmd-cartesian-intro}

\item[\texttt{(ti)}] \texttt{cmd-tuples-intro}

\item[\texttt{(nth-r)}] \texttt{cmd-nth-reduce}

\item[\texttt{(beta)}]
  \texttt{cmd-vnb-lambda-beta} --- $\beta$-reduce
  \texttt{vnb-lambda($x$, $b$)($v$)} in goal.

\item[\texttt{(ii)}] \texttt{cmd-intersection-intro}

\item[\texttt{(ai \textit{f})}]
  \texttt{cmd-antecedent-inference} --- Decompose assumption \textit{f}.

\item[\texttt{(cut \textit{f})}]
  \texttt{cmd-cut} --- Introduce lemma \textit{f}.

\item[\texttt{(ew \textit{t})}]
  \texttt{cmd-exists-witness} --- Provide witness \textit{t} for
  \texttt{forsome} goal.

\item[\texttt{(bc \textit{f})}]
  \texttt{cmd-backchain} --- Backchain on implication \textit{f}.

\item[\texttt{(wk \textit{f})}]
  \texttt{cmd-weaken} --- Remove \textit{f} from assumptions.

\item[\texttt{(ui \textit{k})}]
  \texttt{cmd-union-intro} --- Prove \texttt{x in union(\ldots)} via the
  $k$-th branch.

\item[\texttt{(ue \textit{f})}]
  \texttt{cmd-union-elim} --- Case-split on \texttt{x in union(\ldots)}
  assumption \textit{f}.

\item[\texttt{(ce \textit{f k})}]
  \texttt{cmd-cartesian-elim} --- Extract $k$-th component from Cartesian
  assumption \textit{f}.

\item[\texttt{(te \textit{f k})}]
  \texttt{cmd-tuples-elim} --- Extract $k$-th component from
  \texttt{tuples} assumption \textit{f}.

\item[\texttt{(ie \textit{f k})}]
  \texttt{cmd-intersection-elim} --- Extract $k$-th component from
  intersection assumption \textit{f}.

\item[\texttt{(tfi)}] \texttt{cmd-tfi}

\item[\texttt{(tfi3)}] \texttt{cmd-tfi3}

\item[\texttt{(qed \textit{name})}]
  \texttt{cmd-qed} --- Install the completed proof as a named theorem and
  save the proof script under the same name.

\item[\texttt{(save-proof \textit{name})}]
  Save the current proof script under \textit{name}.

\item[\texttt{(replay-proof \textit{name} [\textit{subst}])}]
  Re-run a saved proof script, optionally with a schematic substitution.

\item[\texttt{(ta \textit{n})}]
  \texttt{cmd-theorem-assumption} --- Add named theorem \textit{n} to
  assumptions.

\item[\texttt{(mac \textit{n})}]
  \texttt{cmd-apply-macete} --- Apply named macete \textit{n} to goal.

\item[\texttt{(inst \textit{f t})}]
  \texttt{cmd-instantiate} --- Instantiate \texttt{forall} assumption
  \textit{f} with term \textit{t}.

\item[\texttt{(focus \textit{n})}]
  \texttt{focus-on} --- Switch focus to the $n$-th open goal (1-based).

\item[\texttt{(show)}]
  \texttt{print-proof-state} --- Display current proof state; updates
  Emacs state buffer.

\item[\texttt{(pp \textit{w})}]
  \texttt{wff->string} --- Display a \texttt{wff} in string syntax.

\item[\texttt{(fa \textit{bs b})}]
  Expand bindings \textit{bs} into nested \texttt{FORALL} formula
  (Section~\ref{sec:multiquant}).

\item[\texttt{(fs \textit{bs b})}]
  Expand bindings \textit{bs} into nested \texttt{FORSOME} formula
  (Section~\ref{sec:multiquant}).

\item[\texttt{(vnb-unwind \textit{expr})}]
  Expand a multi-binding surface expression to nested single-binding form
  (Section~\ref{sec:multiquant}).

\item[\texttt{(vnb-wind \textit{expr})}]
  Collapse a nested single-binding chain into multi-binding form
  (Section~\ref{sec:multiquant}).

\item[\texttt{(parse-string \textit{str})}]
  Parse a string to a raw S-expression.  Note: \texttt{ai}, \texttt{bc},
  \texttt{inst}, \texttt{cut}, \texttt{wk}, \texttt{ue}, \texttt{ce},
  \texttt{ie} accept formula strings directly; \texttt{parse-string} is
  needed only for the low-level \texttt{cmd-*} API.

\item[\texttt{(make-wff-from-string \textit{str})}]
  Parse a string and validate as a \texttt{wff}; use for \texttt{sp}.

\end{description}

Each command (except \texttt{sp}, \texttt{show}, \texttt{pp}, \texttt{fa},
\texttt{fs}, \texttt{vnb-unwind}, \texttt{vnb-wind}, \texttt{parse-string},
\texttt{make-wff-from-string})
mutates the global \texttt{*ps*} and calls \texttt{(show)} automatically.

%% =========================================================
\section{Focus management}

After each command the focus advances automatically to the first
ungrounded subgoal just introduced.  If there are no new subgoals, it
moves to any remaining open goal.

\begin{remark}
Every proof command operates exclusively on the focus node; the other
open goals are untouched.  There can be several open goals simultaneously
(for example, after \texttt{(ci)} on a conjunction goal, or after
\texttt{(cut f)} which always creates two subgoals).  \texttt{(show)}
displays them all, numbering the non-focus goals so you can use those
numbers with \texttt{(focus $n$)}.  Until you call \texttt{(focus $n$)},
the system always works on whichever goal it chose automatically.
\end{remark}

In the interactive session use:
\begin{VNBcode}
(show)       ; list all open goals (numbered) and the current focus
(focus n)    ; switch focus to the n-th open goal (1-based)
\end{VNBcode}

The underlying API (for use outside the interactive session) is:
\begin{VNBcode}
(focus-on ps sqn)            ; focus on a specific sequent node
(focus-on-first-open ps)     ; focus on the first open goal
(proof-open-goals ps)        ; list all open goals
\end{VNBcode}

%% =========================================================
\section{Macetes}

A \emph{macete} is a named, composable rewrite strategy.  When a
theorem is installed, its macete is automatically derived by extracting a
source pattern and replacement pattern from the theorem's logical form:

\begin{center}
\begin{tabular}{ll}
\toprule
Theorem form & Rewrite \\
\midrule
$\forall\bar{x}.\;(s = r)$ & $s \mapsto r$ \\
$\forall\bar{x}.\;(s \Leftrightarrow r)$ & $s \mapsto r$ \\
$\forall\bar{x}.\;(\text{conds} \Rightarrow (s = r))$ & $s \mapsto r$ when conds hold \\
$\forall\bar{x}.\;p$ & $p \mapsto \top$ \\
\bottomrule
\end{tabular}
\end{center}

Macetes can be composed:
\begin{VNBcode}
(macete-series m1 m2 m3)   ; apply in sequence
(macete-repeat m)           ; apply until no change
(macete-try m)              ; apply m, succeed even if m fails
(macete-choice m1 m2 m3)   ; first one that applies
\end{VNBcode}

To install a custom macete:
\begin{VNBcode}
(install-theorem! 'and-comm
  '(FORALL P (FORALL Q (IFF (AND P Q) (AND Q P)))))

; Use it:
(cmd-apply-macete ps 'and-comm)
\end{VNBcode}

%% =========================================================
\section{Example proofs}

\subsection{Identity}

\begin{VNBcode}
(sp (make-wff-from-string "x in NN implies x in NN"))
(di)    ; assume (x in nn), prove (x in nn)
(ass)   ; it is in context
; => Proof complete.
\end{VNBcode}

\subsection{AND elimination}

\begin{VNBcode}
(sp (make-wff-from-string "n in NN and m in NN implies n in NN"))
(di)                          ; assume (n in nn) and (m in nn)
(ai "n in NN and m in NN")   ; split into (n in nn) and (m in nn)
(ass)                         ; prove (n in nn)
\end{VNBcode}

\subsection{VNB axiom instance}

Proving \texttt{"a in b implies a in SET"} from the base axiom:

\begin{VNBcode}
(sp (make-wff-from-string "a in b implies a in SET"))
(di)
(ta 'membership-implies-sethood)
(inst "forall([a, b], a in b implies a in set)" 'a)
(inst "forall([b], a in b implies a in set)" 'b)
(bc "a in b implies a in set")
(ass)
\end{VNBcode}

\subsection{Double negation (classical)}

\begin{VNBcode}
(sp (make-wff-from-string "not(not(x in NN)) implies x in NN"))
(di)                              ; assume not(not((x in nn))), prove (x in nn)
(pbc)                             ; assume not((x in nn)), prove FALSEHOOD
(ai "not(not(x in NN))")         ; not(not(P)) and not(P) yield FALSEHOOD
\end{VNBcode}

\subsection{Ordered tuples and Cartesian products}

\begin{VNBcode}
; From x in A, y in B, z in C prove [x,y,z] in cartesian(A,B,C)
(sp (make-wff-from-string
      "x in A and y in B and z in C implies [x, y, z] in cartesian(A, B, C)"))
(di)
(ai "x in A and y in B and z in C")
(ai "y in B and z in C")
(ci)              ; cartesian-intro generates 3 subgoals
(ass) (ass) (ass)

; NTH projection: nth(2, [x, y, z]) = y
(sp (make-wff-from-string "nth(2, [x, y, z]) = y"))
(nth-r)           ; reduces goal to (y = y)
(rfl)             ; closes by reflexivity
\end{VNBcode}

\subsection{Mathematical structures}

A group can be declared with \texttt{def-structure}.  The carrier is the
underlying set~$G$; the operations are the binary product, the identity
element, and the inverse map.

\begin{VNBcode}
(def-structure 'GROUP
  '(CARRIER)
  '((MUL  (CARTESIAN CARRIER CARRIER) CARRIER) ; function spec
    (UNIT CARRIER)                              ; constant spec: identity element
    (INV  CARRIER CARRIER))                     ; function spec
  '())
; Accessor indices: CARRIER -> 1, MUL -> 2, UNIT -> 3, INV -> 4
\end{VNBcode}

Given \texttt{(IS-GROUP g)} in context, \texttt{cmd-apply-macete} with
\texttt{'IS-GROUP} unfolds to:
\begin{VNBcode}
(AND (IN (CARRIER g) SET)
     (AND (IN (MUL g) (FUN (CARTESIAN (CARRIER g) (CARRIER g))
                            (CARRIER g)))
          (AND (IN (UNIT g) (CARRIER g))
               (IN (INV g) (FUN (CARRIER g) (CARRIER g))))))
\end{VNBcode}

If you prefer to work with raw tuples directly, the n-ary constructors
are still available:
\begin{VNBcode}
; Prove a tuple is in a Cartesian product
(sp (make-wff-from-string
  "G in power(U) and op in fun(cartesian(G,G),G) and e in G and inv in fun(G,G)
   implies [G, op, e, inv] in cartesian(power(U), fun(cartesian(G,G),G), G, fun(G,G))"))
(di)
(ai "G in power(U) and op in fun(cartesian(G,G),G) and e in G and inv in fun(G,G)")
(ai "op in fun(cartesian(G,G),G) and e in G and inv in fun(G,G)")
(ai "e in G and inv in fun(G,G)")
(ci)                              ; cartesian-intro: 4 subgoals
(ass) (ass) (ass) (ass)
\end{VNBcode}

\subsection{Ground arithmetic with \texttt{(arith)}}

The \texttt{(arith)} command (primitive \texttt{cmd-arith}) closes any
goal whose formula reduces to true under exact arithmetic evaluation:

\begin{VNBcode}
(sp (make-wff-from-string "1/3 + 1/6 = 1/2"))
(arith)                              ; closes by exact rational arithmetic

(sp (make-wff-from-string "2 ^ 10 = 1024"))
(arith)                              ; 2^10 evaluates to 1024

(sp (make-wff-from-string "3 + 4 <= 2 * 5"))
(arith)                              ; 7 <= 10
\end{VNBcode}

\subsection{$\beta$-reduction with \texttt{(beta)}}

Applied \texttt{VNB-LAMBDA} terms are reduced in-place, then the
underlying numeric or logical content is closed by another rule:

\begin{VNBcode}
; VNB-LAMBDA has no string syntax; use the S-expression form
; ((VNB-LAMBDA x (* 2 x)) 5) = 10
(sp (make-wff '(= ((VNB-LAMBDA x (* 2 x)) 5) 10)))
(beta)   ; goal becomes ((2 * 5) = 10)
(arith)  ; closes
\end{VNBcode}

\subsection{Auto-installing arithmetic theorems}

\texttt{vnb-create-num-theorem} verifies a closed arithmetic wff by
ground evaluation and installs it as a named theorem.  The theorem is
then available via \texttt{(ta \textit{name})}:

\begin{VNBcode}
(vnb-create-num-theorem 'pythag-3-4-5 "3 ^ 2 + 4 ^ 2 = 5 ^ 2")
; => pythag-3-4-5

(sp (make-wff-from-string "TRUTH implies 3 ^ 2 + 4 ^ 2 = 5 ^ 2"))
(di) (ta 'pythag-3-4-5) (ass)
; => Proof complete.
\end{VNBcode}

False or undecidable formulas are rejected:

\begin{VNBcode}
(vnb-create-num-theorem 'bogus "1 + 1 = 3")
; ERROR: formula evaluates to FALSE

(vnb-create-num-theorem 'bogus2 "x = 5")
; ERROR: not a decidable closed arithmetic sentence
\end{VNBcode}

\subsection{Installing a proven theorem}
\label{sec:qed-and-subst}

A completed proof does not by itself add anything to the theory.  Use
\texttt{(qed \textit{name})} to install the root assertion as a named
theorem:

\begin{VNBcode}
(sp (make-wff-from-string "n in NN and m in NN implies m in NN and n in NN"))
(di) (ai "n in NN and m in NN") (di) (ass) (ass)
(qed 'and-comm)
; => and-comm
; (lookup-theorem 'and-comm) returns the formula S-expression.
\end{VNBcode}

\paragraph{What \texttt{subst-wff} is and is not.}
\texttt{subst-wff} applies capture-avoiding substitution to a wff's
underlying S-expression and returns a new \texttt{wff} (re-validated
through \texttt{make-wff}).  It is \emph{not} a logical inference;
after calling it you must reprove the result.  It is useful for
generating related goals from a known formula:

\begin{VNBcode}
(define w
  (subst-wff '((n . (in x nn)) (m . (in y rr)))
             (lookup-theorem 'and-comm)))
; w is the wff for (x in nn) and (y in rr) implies (y in rr) and (x in nn)
\end{VNBcode}

\subsection{Saving and replaying proof scripts}
\label{sec:replay}

Every interactive command you issue between \texttt{(sp\ \ldots)} and
\texttt{(qed\ \ldots)} is appended to a global accumulator
\texttt{*proof-script*}.  \texttt{(qed \textit{name})} installs the
proven theorem and saves the script under the same name; you can also
call \texttt{(save-proof \textit{name})} on demand.

\begin{VNBcode}
(sp (make-wff-from-string "n in NN and m in NN implies m in NN and n in NN"))
(di) (ai "n in NN and m in NN") (di) (ass) (ass)
(qed 'and-comm)
; theorem AND script installed under the same name

(lookup-proof 'and-comm)
; => ((di) (ai (and (in n nn) (in m nn))) (di) (ass) (ass))
\end{VNBcode}

\paragraph{Replay verbatim.}
\texttt{(replay-proof \textit{name})} re-runs the saved commands
sequentially against the current proof state.  Useful for redoing a
proof in a fresh session, or applying an identical script to a sequent
of the same shape.

%% =========================================================
\chapter{Syntax contexts and local contexts}
\label{sec:syntax-contexts}

\emph{The core context machinery (\texttt{make-wff}, \texttt{wff-free-vars},
\texttt{declare-local-context}, \texttt{undeclare-local-context},
\texttt{get-context}, \texttt{unravel}) is implemented in
\texttt{contexts.scm}.  The bounded form
\texttt{(FORALL (IN $x$\ $A$) $P$)} (and the corresponding
\texttt{FORSOME}) is now expanded automatically by \texttt{make-wff}
to \texttt{(FORALL $x$\ (IMPLIES (IN $x$\ $A$) $P$))}; tuple-destructuring
\texttt{(FORALL (IN (LIST \ldots) \ldots) $P$)} and the parser/printer
expansion to \texttt{NTH} projections remain planned.}

\subsection{Destructuring quantification}

The planned syntax extends the bounded-variable binding form
\texttt{(IN $x$\ $A$)} to allow the bound variable to be a
\emph{tuple pattern}:
\[
  \texttt{(IN (LIST }n_1 \;\ldots\; n_k\texttt{)}\;\; A\texttt{)}
\]
Here $A$ is \emph{any class} --- a Cartesian product, a number system,
a separation class, a structure predicate, or any user-defined class.
There is no requirement that $A$ be a mathematical structure in the sense
of \texttt{def-structure}.

The binding introduces an anonymous implicit tuple variable~$t_0$, adds
\texttt{(IN $t_0$\ $A$)} to the assumptions, and binds each name $n_i$
to the $i$-th component of~$t_0$:
\[
  n_i \;\longmapsto\; \texttt{(NTH }i\;\; t_0\texttt{)}
\]
The variable~$t_0$ never appears to the user.  The scope of
$n_1, \ldots, n_k$ is \emph{lexical}: they have no meaning outside the
quantifier body.

\paragraph{Examples.}

Ranging over a Cartesian product (no structure involved):
\begin{VNBsnip}
(FORALL (IN (LIST x y) (CARTESIAN RR RR))
  (= (+ x y) (+ y x)))

; Expands to:
(FORALL t_0
  (IMPLIES (IN t_0 (CARTESIAN RR RR))
    (= (+ (NTH 1 t_0) (NTH 2 t_0))
       (+ (NTH 2 t_0) (NTH 1 t_0)))))
\end{VNBsnip}

Ranging over a named structure (the structure case is not special):
\begin{VNBsnip}
(FORALL (IN (LIST A + * f) BANACHSPACE)
  (IMPLIES (IN v A) (is-continuous f A RR)))

; Expands to:
(FORALL t_0
  (IMPLIES (IN t_0 BANACHSPACE)
    (IMPLIES (IN v (NTH 1 t_0))
      (is-continuous (NTH 4 t_0) (NTH 1 t_0) RR))))
\end{VNBsnip}

Destructuring bindings compose freely with plain and bounded bindings
(Section~\ref{sec:multiquant}):
\begin{VNBsnip}
(FORALL ((IN (LIST A + * f) BANACHSPACE) (IN v A))
  (is-continuous f A RR))
\end{VNBsnip}

\subsection{Formulas as typed objects: \texttt{wff}}

A \emph{well-formed formula} in this system is not a bare S-expression
but a typed object that records the theory and context in which it was
formed.  In the REPL it prints as:
\begin{VNBsnip}
#[wff N "formula string"]
\end{VNBsnip}
where \texttt{N} is an internal identifier and the quoted string is the
pretty-printed formula.  Every formula in the system is a \texttt{wff},
and its context (theory name) travels with it invisibly.

\texttt{(make-wff expr)} constructs such an object in the current
context.  \texttt{(free-vars p)} returns the variables that are free
\emph{given} that context --- names bound by the current local context
do not count as free.

Two procedures expose the internals of a \texttt{wff}:
\begin{VNBsnip}
(wff->sexp p)        ; => raw S-expression (no output)
(display-contents p) ; prints kind, theory, formula (and context if any)
\end{VNBsnip}
\texttt{wff->sexp} is defined in \texttt{wff.scm};
\texttt{display-contents} is defined in \texttt{contexts.scm}.
Both are available in interactive sessions.

Example (assuming the current theory is \texttt{normal-math}, no local
context active):
\begin{VNBsnip}
(define p (make-wff "forall([x in A], 0 <= funny-name)"))
; => #[wff 12 "forall([x in a], 0 <= funny-name)"]

(free-vars p)
; => (a funny-name)   ; x is bound by the forall
\end{VNBsnip}

\subsection{Local context declarations}

\texttt{declare-local-context} conjures a named instance of a class or
structure into the current session.  The binding
\texttt{(IN (LIST $n_1\ldots n_k$) $A$)} is given a name and installed
as an extension of the current theory:

\begin{VNBsnip}
(declare-local-context (IN (LIST A + * f) BANACHSPACE) bongo)
; => (CONTEXT bongo)
\end{VNBsnip}

Within the scope of \texttt{bongo}, \texttt{A}, \texttt{+}, \texttt{*},
\texttt{f} are no longer free: they are bound to the components of the
anonymous Banach space asserted by \texttt{bongo}.  Any formula
subsequently formed with \texttt{make-wff} carries the extended context:

\begin{VNBsnip}
(define p (make-wff "forall([x in A], 0 <= funny-name)"))
; => #[wff 12 "forall([x in a], 0 <= funny-name)"]   ; tagged with context bongo

(free-vars p)
; => ()    ; A and funny-name now bound by bongo
\end{VNBsnip}

\texttt{(get-context bongo)} returns the full context name
\texttt{(normal-math bongo)}.

\texttt{(undeclare-local-context bongo)} leaves the interactive scope of
\texttt{bongo} but does \emph{not} destroy it.  A formula \texttt{p}
formed inside \texttt{bongo} retains its context tag indefinitely;
the context \texttt{bongo} is still a valid named object that can be
inspected or re-entered.

\subsection{Local contexts as theory extensions}

A local context declaration is syntactic sugar for
\texttt{extend-theory}.  Declaring

\begin{VNBsnip}
(declare-local-context (IN (LIST A + * f) BANACHSPACE) bongo)
\end{VNBsnip}

is equivalent to creating a new theory \texttt{bongo} that extends
\texttt{normal-math} with the single axiom:
\[
  \texttt{(IN (LIST A + * f) BANACHSPACE)}
\]
This axiom asserts the existence of a Banach space whose components are
exactly \texttt{A}, \texttt{+}, \texttt{*}, \texttt{f}.  The names
cease to be free variables because the theory now witnesses them.
\texttt{undeclare-local-context} simply records that the interactive
session is no longer inside that theory extension; it does not retract
the axiom or invalidate any work done under it.

\subsection{The \texttt{unravel} operation}

A formula tagged with a local context can be \emph{unravelled}: the
context binding is absorbed into the formula itself as an explicit
destructuring quantifier, producing a self-contained statement that no
longer depends on the named context.

\begin{VNBsnip}
; p = #[wff 12 "forall([x in a], 0 <= funny-name)"]   ; tagged with context bongo
; bongo binds: [A, +, *, f] in BANACHSPACE

(unravel p)
; => #[wff 13 "forall([a, +, *, f in banachspace], forall([x in a], 0 <= funny-name))"]
\end{VNBsnip}

Each local context layer is peeled off in order, outermost first, and
replaced by the corresponding destructuring \texttt{FORALL}.  The result
lives in the base theory (\texttt{normal-math}) with no remaining
context dependency.

\subsection{The syntax context}

Both the destructuring quantifier and the local context declaration
operate through a \emph{syntax context}: a table mapping user-facing
names to \texttt{nth($k$, $t_0$)} projections of the implicit tuple
instance.  The syntax context is consulted at two points:

\begin{itemize}[label=\textendash, leftmargin=1cm]
\item \textbf{Parser (input):} names in the syntax context are expanded
  to their \texttt{nth} forms before a formula enters the proof state.
\item \textbf{Printer (output):} \texttt{nth} projections of a known
  implicit instance are collapsed back to the declared names when the
  proof state is displayed.
\end{itemize}

The syntax context must travel with the proof state and with
\texttt{wff} objects.  Without the inverse mapping at display
time, the user would see anonymous \texttt{NTH} projections instead of
the declared names, defeating the purpose of the mechanism.

The two scoping disciplines share the same underlying syntax context
machinery:

\begin{center}
\begin{tabular}{lll}
\toprule
Form & Scope & Basis \\
\midrule
Destructuring quantifier &
  Lexical: quantifier body &
  Formula structure \\
\texttt{declare-local-context} &
  Dynamic: until \texttt{undeclare} &
  Named theory extension \\
\bottomrule
\end{tabular}
\end{center}

\subsection{Parser and printer}
\label{sec:parser-printer}

Four entry points provide string-level access to VNB formulas:

\begin{VNBsnip}
(make-wff-from-string str)  ; string -> <wff>  (validates via make-wff)
(parse-string str)          ; string -> raw S-expression (no validation)
(wff->string w)             ; <wff>  -> string
(pp w)                      ; display wff->string then newline
\end{VNBsnip}

\noindent
\texttt{parse-string} converts a string to a raw S-expression without
validation; it is used internally by the interactive commands and is
available for low-level \texttt{cmd-*} calls.  For normal use, pass
formula strings directly to \texttt{ai}, \texttt{bc}, \texttt{inst},
\texttt{cut}, \texttt{wk}, and \texttt{vnb-create-num-theorem}.
\texttt{make-wff-from-string} is the right choice when constructing a
formula to prove (for \texttt{sp}).

\paragraph{Operator precedence (high to low).}
\begin{center}
\small
\begin{tabular}{lll}
\toprule
Level & Operators & Associativity \\
\midrule
primary & atoms, \texttt{f(a,b)}, \texttt{[a,b]}, \texttt{(expr)} & --- \\
power & \texttt{\textasciicircum} & right \\
unary & \texttt{-x} & right \\
mul & \texttt{*}, \texttt{/} & left (n-ary \texttt{*}) \\
add & \texttt{+}, \texttt{-} & left (n-ary \texttt{+}) \\
cmp & \texttt{in}, \texttt{=}, \texttt{==}, \texttt{<=} & binary \\
and & \texttt{and} & left (n-ary) \\
or & \texttt{or} & left (n-ary) \\
iff & \texttt{iff} & left (n-ary) \\
implies & \texttt{implies} & right \\
\bottomrule
\end{tabular}
\end{center}

\paragraph{Division and negation sugar.}
\begin{VNBsnip}
x / y    =>  (* x (recip y))   ; division is syntactic sugar
-x       =>  (- x)             ; unary minus; lower prec than ^
-x ^ 2   =>  (- (power x 2))  ; i.e. -(x^2), not (-x)^2
\end{VNBsnip}

\paragraph{Quantifier syntax.}
Quantifiers use the binding-list form:
\begin{VNBsnip}
forall([x], body)              ; bare variable (all classes)
forall([x in A], body)         ; guarded variable
forall([[a, b, c] in A], body) ; tuple destructuring
forall([x in A, y, z in B], body)  ; mixed bindings
forsome([x in NN], body)       ; existential; same forms
\end{VNBsnip}

\noindent
The bracket syntax after the function name distinguishes quantifiers
from ordinary function calls.

\paragraph{Display conventions (\texttt{wff->string} / \texttt{pp}).}
\begin{itemize}[label=\textendash, leftmargin=1cm]
\item \textbf{Quantifiers} print in binding-list form:
  the primitive expanded form is reverse-engineered.
  \texttt{forall([x in nn, y], p(x, y))}
\item \textbf{Logical connectives} \texttt{and}, \texttt{or}, \texttt{iff}
  print n-ary infix with surrounding parentheses:
  \texttt{((p and q) and r)}.
  \texttt{implies} prints binary infix: \texttt{(p implies q)}.
\item \textbf{Comparisons} \texttt{in}, \texttt{=}, \texttt{==},
  \texttt{<=} print binary infix: \texttt{(x in nn)}, \texttt{(x = y)}.
\item \textbf{Arithmetic} \texttt{+}, \texttt{*} print n-ary infix:
  \texttt{(x + y + z)}.  Binary \texttt{-} prints as \texttt{(x - y)};
  unary as \texttt{-x}.  Exponentiation as \texttt{(x \textasciicircum{} y)}.
\item \textbf{Lists} \texttt{[a, b, c]} in both input and output: the
  list syntax is the same in string form (internally stored as \texttt{LIST}).
\item \textbf{All other compound terms} use functional notation
  \texttt{f(x, y, z)}: \texttt{nth}, \texttt{cartesian}, \texttt{fun},
  \texttt{power} (unary set power), \texttt{not}, and user-defined heads.
  Nullary symbols print without parentheses.
\end{itemize}

\paragraph{Example round-trips.}
\begin{VNBsnip}
(pp (make-wff-from-string "forall([x in NN, y in NN], x + y = y + x)"))
; => forall([x in nn, y in nn], ((x + y) = (y + x)))

(pp (make-wff-from-string "not(P) implies FALSEHOOD"))
; => (not(p) implies falsity)

(pp (make-wff-from-string "2 ^ 10 = 1024"))
; => ((2 ^ 10) = 1024)

(expression->string '(forall x (implies (in x nn)
                       (forall y (implies (in y nn) (= (+ x y) (+ y x)))))))
; => "forall([x in nn, y in nn], ((x + y) = (y + x)))"
\end{VNBsnip}

\noindent
All identifiers are case-folded to lowercase by MIT~Scheme, so
\texttt{NN} and \texttt{nn} are the same symbol; \texttt{FORALL} and
\texttt{forall} are also identical.  The printer always emits lowercase.

%% =========================================================
\chapter{Source file overview}

\begin{description}

\item[\texttt{expressions.scm}]
  Expression representation, free variables, substitution
  (capture-avoiding), alpha-equivalence.

\item[\texttt{sequents.scm}]
  Sequents (list of assumptions + assertion), context operations.

\item[\texttt{deduction-graphs.scm}]
  Sequent nodes, inference nodes, the deduction graph record type,
  grounding propagation.

\item[\texttt{primitive-inferences.scm}]
  The kernel: all proof rules as Scheme procedures.  The only code
  allowed to extend the deduction graph.

\item[\texttt{macetes.scm}]
  Pattern matching, elementary macete construction from theorems,
  compound macete combinators.

\item[\texttt{theory.scm}]
  Theory record, axiom/theorem registration, the VNB base theory.

\item[\texttt{structures.scm}]
  \texttt{def-structure}: registers structure definitions, installs
  NTH-based accessor macetes, auto-generates \texttt{IS-\textit{NAME}}
  axioms.

\item[\texttt{number-systems.scm}]
  Declares \texttt{NN}, \texttt{ZZ}, \texttt{QQ}, \texttt{RR},
  \texttt{CC} as constants and installs their arithmetic axioms
  (sethood, inclusion chain, ring/field laws, absolute value,
  conjugation).

\item[\texttt{wff.scm}]
  \texttt{<wff>} record type (\texttt{formula}, \texttt{theory},
  \texttt{contexts} fields); \texttt{wff?}, \texttt{wff-formula},
  \texttt{wff-theory}; \texttt{wff-child} (inherit theory from parent),
  \texttt{wff-in-theory} (bare constructor), \texttt{wff-equiv?}.

\item[\texttt{contexts.scm}]
  \texttt{make-wff} (snapshot current theory and context stack),
  \texttt{wff-free-vars}, \texttt{declare-local-context},
  \texttt{undeclare-local-context}, \texttt{get-context},
  \texttt{unravel}, \texttt{print-wff}, \texttt{display-contents}
  (browse a wff's formula, theory, and kind).

\item[\texttt{arith-eval.scm}]
  Ground arithmetic decision procedure.  Evaluates closed formulas over
  $\mathit{NN}/\mathit{ZZ}/\mathit{QQ}/\mathit{RR}/\mathit{CC}$ using
  Scheme's exact arithmetic.  Handles \texttt{+}, \texttt{*},
  \texttt{-}, \texttt{recip}, \texttt{abs}, \texttt{conjugate},
  \texttt{succ}, \texttt{<=}, \texttt{=}, \texttt{IN}, and
  \texttt{power}.  Provides \texttt{vnb-do-arith} (standalone
  predicate), \texttt{vnb-create-num-theorem} (verify and install a
  numeric theorem), and \texttt{pi-arith!} / \texttt{cmd-arith} (closes
  proof goals).

\item[\texttt{proof-commands.scm}]
  High-level command wrappers with automatic focus management.  Also
  \texttt{cmd-qed} (install a completed proof as a theorem) and
  \texttt{subst-wff} (meta-level substitution into a wff).

\item[\texttt{interactive.scm}]
  Interactive short-form commands (\texttt{sp}, \texttt{di},
  \texttt{ai}, \ldots), \texttt{fa}/\texttt{fs} quantification helpers,
  and the \texttt{show} display function.  Loaded at session startup.

\item[\texttt{vnb.el}]
  GNU Emacs / XEmacs interface.  REPL buffer (\texttt{*VNB*}),
  proof-state display buffer (\texttt{*VNB State*}), command scratch
  buffer (\texttt{*VNB Commands*}), command alist and completion,
  \texttt{vnb-eval-string} for programmatic evaluation.

\end{description}

%% =========================================================

\begin{thebibliography}{9}
\bibitem{IMPS}
W.~M.~Farmer, J.~D.~Guttman, F.~J.~Thayer,
\textit{IMPS: An Interactive Mathematical Proof System},
Journal of Automated Reasoning 11(2), 1993.
\end{thebibliography}

\end{document}
